% !TEX program = pdflatex
% =====================================================================
%  SAPHO — Manual do Usuário
%  Documentação viva da plataforma SAPHO (AURORA + YANC + toolchain).
%  Compilar:  pdflatex main  (3x, para sumário e referências)
%
%  Projeto gráfico: livro clássico de ensino de tecnologia.
%  Tipografia Palatino (mathpazo, algarismos antigos), paleta sóbria
%  de tinta e bronze, parágrafos com recuo, versalete nos elementos
%  estruturais. Estilo de escrita: CONVENCOES-GERAIS do grupo (sem
%  travessões, sem negrito no corpo, prosa encadeada).
% =====================================================================
\documentclass[11pt,a4paper,twoside,openright]{book}

\usepackage[utf8]{inputenc}
\usepackage[T1]{fontenc}
\usepackage[brazil]{babel}

% ---- tipografia ----
\usepackage{lmodern}                 % base (fornece a mono lmtt)
\usepackage[osf,sc]{mathpazo}        % Palatino com algarismos antigos
\usepackage[scale=0.90]{tgheros}     % sans discreta p/ elementos de interface
\renewcommand{\ttdefault}{lmtt}      % mono clara da Latin Modern
\linespread{1.10}
\usepackage{textcomp}
\usepackage[letterspace=60]{microtype}

% ---- página ----
\usepackage[a4paper,
  inner=3.0cm, outer=3.6cm,
  top=3.4cm, bottom=3.8cm,
  headsep=0.9cm]{geometry}
\setlength{\parindent}{1.25em}
\setlength{\parskip}{1pt plus 1pt}

\usepackage{xcolor}
\usepackage{graphicx}
\usepackage{booktabs}
\usepackage{array}
\usepackage{tabularx}
\usepackage{longtable}
\usepackage{enumitem}
\usepackage{multicol}
\usepackage{csquotes}
\usepackage{listings}
\usepackage{titlesec}
\usepackage{framed}
\usepackage{fancyhdr}
\setlength{\headheight}{14pt}
\usepackage[hidelinks]{hyperref}
\usepackage{url}

\setlist{itemsep=2pt,topsep=6pt,parsep=0pt}
\setlength{\tabcolsep}{6pt}
\renewcommand{\arraystretch}{1.35}

% ---- paleta: tinta e bronze ----
\definecolor{ink}{HTML}{23282F}        % quase-preto dos títulos
\definecolor{bronze}{HTML}{8A6D3B}     % ornamentos, números de capítulo
\definecolor{slate}{HTML}{44546A}      % links, rótulo de informação
\definecolor{wine}{HTML}{7B3B3B}       % rótulo de cuidado
\definecolor{amber}{HTML}{95691D}      % rótulo de aviso
\definecolor{sage}{HTML}{4E6E58}       % rótulo de nota
\definecolor{cinza}{HTML}{707070}      % fólios e cabeços
\definecolor{regua}{HTML}{C9C2B4}      % réguas finas
\definecolor{codeKw}{HTML}{35506B}     % palavras-chave de código
\definecolor{codeCom}{HTML}{7A7A6E}    % comentários
\definecolor{codeStr}{HTML}{7B3B3B}    % strings

\hypersetup{
  pdftitle={SAPHO --- Manual do Usuário},
  pdfauthor={NIPS-CERN / UFJF},
  colorlinks=true, linkcolor=slate, urlcolor=slate, citecolor=slate}

% ---- colunas tabularx ----
\newcolumntype{Y}{>{\raggedright\arraybackslash}X}

% ---- código: sem moldura, régua fina à esquerda ----
\lstset{
  basicstyle=\ttfamily\small,
  keywordstyle=\color{codeKw}\bfseries,
  commentstyle=\color{codeCom}\itshape,
  stringstyle=\color{codeStr},
  frame=leftline, rulecolor=\color{regua}, framerule=0.8pt, framesep=10pt,
  xleftmargin=14pt,
  breaklines=true, columns=fullflexible, keepspaces=true,
  showstringspaces=false, tabsize=2, upquote=true, extendedchars=true,
  aboveskip=14pt, belowskip=14pt,
  literate=%
    {á}{{\'a}}1 {à}{{\`a}}1 {â}{{\^a}}1 {ã}{{\~a}}1
    {é}{{\'e}}1 {ê}{{\^e}}1 {í}{{\'i}}1
    {ó}{{\'o}}1 {ô}{{\^o}}1 {õ}{{\~o}}1
    {ú}{{\'u}}1 {ü}{{\"u}}1 {ç}{{\c c}}1
    {Á}{{\'A}}1 {É}{{\'E}}1 {Ç}{{\c C}}1 {Ø}{{\O}}1 {ø}{{\o}}1
    {±}{{$\pm$}}1 {→}{{$\rightarrow$}}1 {≈}{{$\approx$}}1
}
\lstdefinelanguage{CMM}{morekeywords={void,int,float,comp,if,else,while,for,do,switch,case,
  default,break,continue,return,
  in,out,fin,fout,abs,sign,norm,pset,copy,sqrt,atan,sin,cos,tan,sinh,cosh,tanh,exp,log,pow,
  floor,ceil,round,real,imag,fase,mod2,complex,conj},
  morecomment=[l]{//},morecomment=[s]{/*}{*/},morestring=[b]",
  morecomment=[l][\color{codeKw}\bfseries]{\#}}
\lstdefinestyle{cmm}{language=CMM}
\lstdefinelanguage{Verilog}{morekeywords={module,endmodule,input,output,inout,wire,reg,
  assign,always,posedge,negedge,begin,end,if,else,initial,parameter,localparam,integer},
  morecomment=[l]{//},morecomment=[s]{/*}{*/},morestring=[b]"}

% ---- títulos: capítulo com número em bronze, versalete no rótulo ----
\titleformat{\chapter}[display]
  {\normalfont}
  {\filright\textls[120]{\scshape\color{cinza}\MakeLowercase{capítulo}}~{\fontsize{42}{42}\selectfont\color{bronze}\thechapter}}
  {14pt}
  {\filright\huge\bfseries\color{ink}}
  [\vspace{6pt}{\color{regua}\titlerule[0.8pt]}]
\titlespacing*{\chapter}{0pt}{22pt}{34pt}

\titleformat{\section}
  {\Large\bfseries\color{ink}}{\textcolor{bronze}{\thesection}}{0.7em}{}
\titlespacing*{\section}{0pt}{26pt}{10pt}
\titleformat{\subsection}
  {\large\bfseries\color{ink}}{\textcolor{bronze}{\thesubsection}}{0.7em}{}
\titlespacing*{\subsection}{0pt}{18pt}{7pt}
\titleformat{\subsubsection}
  {\normalsize\bfseries\itshape\color{ink}}{\thesubsubsection}{0.7em}{}
\titlespacing*{\subsubsection}{0pt}{14pt}{5pt}

% partes
\titleformat{\part}[display]
  {\centering\normalfont}
  {\textls[160]{\scshape\color{cinza}\MakeLowercase{parte}}~{\Huge\color{bronze}\thepart}}
  {24pt}
  {\Huge\bfseries\color{ink}}
\titlespacing*{\part}{0pt}{6cm}{0pt}

% ---- cabeços: versalete espaçado, fólio externo, sem régua ----
\pagestyle{fancy}
\fancyhf{}
\fancyhead[LE]{\small\color{cinza}\thepage\qquad\textls[80]{\scshape\MakeLowercase{\leftmark}}}
\fancyhead[RO]{\small\color{cinza}\textls[80]{\scshape\MakeLowercase{\rightmark}}\qquad\thepage}
\renewcommand{\headrulewidth}{0pt}
\renewcommand{\chaptermark}[1]{\markboth{#1}{}}
\renewcommand{\sectionmark}[1]{\markright{#1}}
\fancypagestyle{plain}{\fancyhf{}\fancyfoot[C]{\small\color{cinza}\thepage}\renewcommand{\headrulewidth}{0pt}}

% ---- macros ----
\newcommand{\cmm}{C\textpm{}}
\newcommand{\repo}[1]{\texttt{nipscernlab/#1}}
\newcommand{\file}[1]{\path{#1}}
\newcommand{\tool}[1]{\textsc{#1}}
\newcommand{\code}[1]{\lstinline[basicstyle=\ttfamily]|#1|}
\newcommand{\botao}[1]{\textsf{#1}}                     % botão da interface (sans, sem negrito)
\newcommand{\menu}[1]{\textsf{#1}}                      % item de menu
\newcommand{\tecla}[1]{\fbox{\footnotesize\textsf{#1}}} % tecla de atalho

% ---- caixas: régua fina à esquerda + rótulo em versalete colorido ----
\newenvironment{caixa}[2]{%
  \def\FrameCommand{{\color{#1}\vrule width 0.8pt}\hspace{12pt}}%
  \MakeFramed{\advance\hsize-\width\FrameRestore}%
  \vspace{2pt}\noindent\textcolor{#1}{\textls[100]{\scshape #2}}\\[3pt]\small\noindent\ignorespaces}%
 {\endMakeFramed}
\newenvironment{nota}{\begin{caixa}{sage}{nota}}{\end{caixa}}
\newenvironment{info}{\begin{caixa}{slate}{informação}}{\end{caixa}}
\newenvironment{aviso}{\begin{caixa}{amber}{aviso}}{\end{caixa}}
\newenvironment{perigo}{\begin{caixa}{wine}{cuidado}}{\end{caixa}}

% ---- marcador de figura pendente (print a ser tirado) ----
\newcounter{figpend}
\newcommand{\figurapendente}[2]{%
  \refstepcounter{figpend}%
  \par\vspace{16pt}
  \begin{center}
  \begin{minipage}{0.86\linewidth}
  \centering
  {\color{regua}\rule{\linewidth}{0.8pt}}\\[8pt]
  {\small\textcolor{bronze}{\textls[100]{\scshape figura \thefigpend}} \,\textit{#1}}\\[4pt]
  {\footnotesize\color{cinza}#2}\\[8pt]
  {\color{regua}\rule{\linewidth}{0.8pt}}
  \end{minipage}
  \end{center}
  \vspace{16pt}}

\begin{document}
\frontmatter

% ===================== capa =====================
\begin{titlepage}
\centering
\vspace*{2.8cm}
{\color{regua}\rule{0.35\linewidth}{0.8pt}}\\[26pt]
{\fontsize{58}{58}\selectfont\textls[180]{\scshape sapho}\par}
\vspace{16pt}
{\Large\color{ink}\textit{Manual do Usuário}\par}
\vspace{34pt}
{\normalsize\color{cinza} Concepção, programação, simulação e visualização\\ de \textit{soft-processors} com a IDE AURORA\par}
\vfill
{\small\textls[80]{\scshape\color{ink} nips-cern}\\[4pt]
\small\color{cinza} Núcleo de Instrumentação e Processamento de Sinais\\
Universidade Federal de Juiz de Fora\par}
\vspace{18pt}
{\small\color{cinza}\url{https://nipscern.com/sapho}\par}
\vspace{26pt}
{\color{regua}\rule{0.35\linewidth}{0.8pt}}\\[14pt]
{\footnotesize\color{cinza}\textit{Documento vivo. Acompanha as versões da plataforma.}\par}
\vspace*{1.6cm}
\end{titlepage}

\tableofcontents

% !TEX root = ../main.tex
\chapter{Prefácio}
\label{cap:prefacio}

Este livro é o manual de uso da plataforma \tool{SAPHO}: da instalação ao primeiro processador funcionando em simulação, e de lá até os recursos mais avançados da IDE \tool{AURORA} e da linguagem \cmm{}. Ele foi escrito para quem \textbf{nunca viu o SAPHO} --- nenhum conhecimento prévio da plataforma é assumido, embora familiaridade básica com programação em C e com a ideia de circuitos digitais ajude.

\section*{Como este livro está organizado}

O livro é dividido em módulos, um por componente da plataforma --- cada componente do SAPHO é, na prática, um projeto por si só, e recebeu aqui a sua própria mini-documentação:

\begin{description}
  \item[Parte I --- Primeiros passos.] O que é o SAPHO, instalação e um tour pela interface da AURORA.
  \item[Parte II --- Projetos e processadores.] O arquivo de projeto \file{.spf}, a arquitetura do processador SAPHO e o Hub de Processadores.
  \item[Parte III --- A linguagem \cmm{}.] A referência completa da linguagem: estrutura, tipos, operadores, E/S, biblioteca padrão e recursos avançados.
  \item[Parte IV --- O editor.] O editor Monaco e todo o apoio à edição: abas, divisão de painéis, diagnósticos, formatação, busca.
  \item[Parte V --- Compilação e simulação.] Os botões da toolbar, o pipeline do YANC, os seis terminais e os quatro caminhos de simulação.
  \item[Parte VI --- Visualização.] Formas de onda (GTKWave e Surfer) e o visualizador de RTL PRISM.
  \item[Parte VII --- Aurora Intelligence.] O assistente de IA integrado: configuração, provedores, ferramentas e permissões.
  \item[Parte VIII --- Ferramentas de apoio.] Dagr (controle de versão Git) e as configurações da IDE.
  \item[Parte IX --- Manutenção.] O atualizador automático e a solução de problemas.
  \item[Apêndices.] Referências rápidas: diretivas, biblioteca padrão, atalhos de teclado e glossário.
\end{description}

Um exemplo único costura o livro: o processador \code{media_movel}, um filtro de média móvel construído do zero no projeto \code{MeuFiltro}. Ele nasce no Capítulo~\ref{cap:criando-processadores}, ganha código na Parte~III, compila na Parte~V e aparece nas ondas e no RTL na Parte~VI.

\section*{Convenções usadas neste livro}

\begin{description}
  \item[\botao{Botão}] Elementos clicáveis da interface: botões, itens de menu de contexto.
  \item[\tecla{Ctrl+S}] Atalhos de teclado.
  \item[\code{codigo}] Nomes de comandos, diretivas, variáveis, valores literais.
  \item[\file{arquivo.cmm}] Nomes de arquivos e caminhos.
  \item[\tool{Nome}] Ferramentas e componentes da plataforma.
\end{description}

As caixas coloridas seguem o vocabulário de cores da própria AURORA:

\begin{nota}
\textbf{Nota} (menta): dicas e informações que enriquecem o assunto ao lado.
\end{nota}

\begin{info}
\textbf{Informação} (ciano): contexto adicional, referências cruzadas, detalhes de bastidores.
\end{info}

\begin{aviso}
\textbf{Aviso} (âmbar): restrições e comportamentos que costumam surpreender.
\end{aviso}

\begin{perigo}
\textbf{Cuidado} (vermelho): ações destrutivas ou erros com consequências sérias.
\end{perigo}

\section*{Um documento vivo}

Este manual acompanha o desenvolvimento da plataforma e é atualizado junto com a AURORA, o YANC e os demais componentes. Cada afirmação sobre comportamento foi extraída do código-fonte dos repositórios oficiais --- não de documentos intermediários --- na data da edição. Se você encontrar divergência entre o manual e o software, o software mudou primeiro: procure a edição mais recente ou abra uma issue nos repositórios do \texttt{nipscernlab}.


% ===================== módulos =====================
\mainmatter

\part{Primeiros passos}
% !TEX root = ../main.tex
\chapter{O que é o SAPHO}
\label{cap:o-que-e-sapho}

\section{A plataforma em uma frase}

\tool{SAPHO} (\emph{Scalable-Architecture Processor for Hardware Optimization}) é uma plataforma completa para criar, programar, simular e visualizar \emph{soft-processors}: você escreve um programa numa linguagem parecida com C, chamada \cmm{} (lê-se \enquote{C mais-menos}), e a plataforma gera um processador dedicado em Verilog, sob medida para aquele programa, pronto para ser sintetizado em FPGA. Tudo acontece dentro de uma única aplicação de desktop, a IDE \tool{AURORA} (\emph{Advanced Utility Running Optimized Resource Architectures}), que roda localmente no Windows, sem depender de nuvem.

\section{O que é um soft-processor}

Um processador convencional (como o do seu computador) é um chip fixo: sua arquitetura foi decidida na fábrica e não muda. Um \emph{soft-processor}, por outro lado, é um processador descrito em uma linguagem de descrição de hardware (HDL, no caso o Verilog) e implementado na malha reconfigurável de um FPGA. Isso significa que a arquitetura --- largura da palavra, conjunto de instruções, quantidade de portas de entrada e saída, tamanho das memórias --- é escolhida por você, e pode ser alterada e recompilada em minutos.

O SAPHO leva essa ideia ao extremo: o processador gerado contém \emph{apenas} o hardware que o seu programa usa. Se o programa não usa divisão, o divisor não existe no circuito final. Se não usa ponto flutuante, nenhuma lógica de ponto flutuante é gerada. O resultado é um processador enxuto, previsível e eficiente, ideal para instrumentação científica e processamento de sinais em tempo real.

\section{As peças da plataforma}
\label{sec:pecas-da-plataforma}

O SAPHO é um guarda-chuva que reúne vários componentes. Cada um deles tem um módulo próprio neste manual.

\begin{longtable}{@{}p{0.24\linewidth}p{0.70\linewidth}@{}}
\toprule
\textbf{Componente} & \textbf{O que é} \\
\midrule
\endhead
\tool{AURORA} & \emph{Advanced Utility Running Optimized Resource Architectures}: a IDE de desktop --- editor de código, gerenciador de projetos, botões de compilação e simulação, terminais, visualizadores. É o programa que você instala e abre. \\
\tool{YANC} & \emph{Yet Another Compiler}: a suíte de compiladores que traduz seu programa \cmm{} (ou C++) no processador Verilog, nas imagens de memória e no \emph{testbench}. Roda por baixo dos panos quando você clica nos botões de compilação. \\
Processador \tool{SAPHO} & O \emph{soft-processor} em si: o circuito Verilog parametrizável gerado pelo YANC (Capítulo~\ref{cap:processador-sapho}). \\
\emph{Toolchain bundle} & Um pacote portátil de ferramentas de EDA de código aberto (Icarus Verilog, Verilator, Yosys, GCC, Python + cocotb) que a AURORA baixa automaticamente na primeira execução. \\
\tool{PRISM} & \emph{Processor Rendering Interface for Schematic Models}: o visualizador de RTL embutido, que desenha o diagrama do circuito sintetizado, navegável módulo a módulo (Capítulo~\ref{cap:prism}). \\
\tool{GTKWave} / \tool{Surfer} & Os dois visualizadores de formas de onda disponíveis para inspecionar simulações (Capítulo~\ref{cap:ondas}). \\
\tool{Aurora Intelligence} & O assistente de IA integrado, que conversa sobre o projeto e age sobre a IDE por ferramentas curadas (Capítulo~\ref{cap:aurora-intelligence}). \\
\bottomrule
\end{longtable}

\section{O fluxo de trabalho, de ponta a ponta}

O caminho típico de quem usa o SAPHO é sempre o mesmo, e este manual o percorre na ordem:

\begin{enumerate}
  \item \textbf{Criar um projeto} --- uma pasta com um arquivo \file{.spf} que descreve tudo (Capítulo~\ref{cap:projetos}).
  \item \textbf{Criar um processador} --- escolher nome e parâmetros de arquitetura (Capítulo~\ref{cap:criando-processadores}).
  \item \textbf{Programar em \cmm{}} --- escrever o comportamento no editor (Parte~III).
  \item \textbf{Compilar} --- o YANC transforma o programa em Verilog + memórias + testbench (Capítulo~\ref{cap:compilacao}).
  \item \textbf{Simular} --- Icarus Verilog ou Verilator executam o circuito (Capítulo~\ref{cap:simulacao}).
  \item \textbf{Visualizar} --- formas de onda no GTKWave/Surfer, circuito no PRISM (Parte~VI).
  \item \textbf{Sintetizar no FPGA} --- levar o \file{.v} e os \file{.mif} gerados para a ferramenta do fabricante (Quartus, Vivado etc.), fora da AURORA.
\end{enumerate}

\section{O exemplo que acompanha este manual}
\label{sec:exemplo-condutor}

Para que cada conceito apareça em contexto, este manual constrói um processador de exemplo do zero e o acompanha até a forma de onda final: o \code{media_movel}, um filtro de média móvel que lê amostras de um sinal por uma porta de entrada, calcula a média das últimas amostras e escreve o resultado numa porta de saída. É um exemplo pequeno, mas exercita o essencial: entrada e saída, aritmética, laço principal e simulação com estímulo externo.

\section{O laboratório por trás: NIPS-CERN}

O SAPHO é desenvolvido e mantido pelo \textbf{NIPS-CERN}, núcleo de pesquisa da \textbf{Universidade Federal de Juiz de Fora (UFJF)}, em Minas Gerais. O laboratório colabora com o CERN no experimento \textbf{ATLAS} do LHC, onde técnicas de processamento de sinais em FPGA --- o habitat natural do SAPHO --- são aplicadas à instrumentação dos calorímetros. A plataforma também é usada no ensino, na disciplina de Dispositivos Lógicos Programáveis (DLP) da UFJF.

Todo o ecossistema é de código aberto, hospedado na organização \texttt{nipscernlab} no GitHub. O canal oficial de download para o usuário final é \url{https://nipscern.com/sapho}.

\begin{nota}
Este é um documento vivo: ele acompanha as versões da plataforma e é atualizado junto com a AURORA, o YANC e os demais componentes. Se algo na sua tela não corresponde ao manual, verifique se a sua instalação está atualizada (Capítulo~\ref{cap:atualizacoes}).
\end{nota}

% !TEX root = ../main.tex
\chapter{Instalação}
\label{cap:instalacao}

\section{Requisitos}

\begin{itemize}
  \item \textbf{Windows 64 bits} (Windows 10 ou 11). A AURORA é, hoje, um aplicativo exclusivamente Windows.
  \item Cerca de 2\,GB de espaço em disco (o aplicativo já traz embutida toda a cadeia de ferramentas de simulação e síntese).
  \item Nenhum outro programa é necessário: compiladores, simuladores e visualizadores vêm no pacote.
\end{itemize}

\section{Baixando e instalando}

\begin{enumerate}
  \item Acesse \url{https://nipscern.com/sapho} e baixe o instalador --- um arquivo chamado \file{sapho-aurora-Setup-vX.Y.Z.exe} (onde \code{X.Y.Z} é a versão).
  \item Dê dois cliques no executável. O instalador é do tipo \emph{one-click}: não há telas de opções nem escolha de pasta --- ele instala, cria os atalhos na área de trabalho e no menu Iniciar e abre o aplicativo.
\end{enumerate}

\figurapendente{Página de download em nipscern.com/sapho}{Mostrar o botão de download do instalador.}

\figurapendente{Primeira abertura da AURORA: tela de splash e tela de boas-vindas}{Capturar a splash com a barra de progresso e, em seguida, a tela de boas-vindas com o fundo animado de aurora boreal.}

Durante a instalação são registrados dois vínculos com o sistema:
\begin{itemize}
  \item a extensão \file{.spf} (arquivo de projeto SAPHO) passa a abrir com a AURORA, com ícone próprio --- dois cliques num \file{.spf} no Explorer abrem o projeto;
  \item o protocolo \code{sapho://} é registrado para integrações futuras via navegador.
\end{itemize}

\begin{nota}
O instalador ainda não é assinado digitalmente; o Windows SmartScreen pode exibir um aviso de \enquote{editor desconhecido} na primeira execução. Clique em \menu{Mais informações $\to$ Executar assim mesmo}.
\end{nota}

\section{O que é instalado, e onde}

\begin{longtable}{@{}p{0.34\linewidth}p{0.60\linewidth}@{}}
\toprule
\textbf{Local} & \textbf{Conteúdo} \\
\midrule
\endhead
Pasta do aplicativo & O executável da AURORA e a pasta \file{components/} com toda a toolchain: compiladores do YANC, simuladores (Icarus Verilog, Verilator), Yosys, GTKWave, Python com cocotb e os moldes de HDL do processador. \\
\file{\%APPDATA\%\textbackslash Aurora-IDE\textbackslash} & Dados do usuário: logs (\file{logs/main.log}), lista de projetos recentes (\file{aurora-recents.json}), registro de versão para o pós-atualização. \\
Cofre de credenciais do Windows & Chaves de API da Aurora Intelligence e token do GitHub, cifrados pelo sistema operacional (DPAPI) --- nunca em texto puro. \\
\bottomrule
\end{longtable}

Os seus projetos ficam onde você quiser: a AURORA pergunta a pasta ao criar cada projeto (Capítulo~\ref{cap:projetos}).

\section{Primeira execução}

Ao abrir, a AURORA exibe uma tela de \emph{splash} com o progresso real de inicialização e, em seguida, a \textbf{tela de boas-vindas}, com o fundo animado de aurora boreal, os botões \botao{Novo Projeto} e \botao{Abrir Projeto} e a lista de projetos \textbf{Recentes} (vazia na primeira vez). A partir daí, siga para o Capítulo~\ref{cap:tour-interface} para conhecer a interface, ou direto para o Capítulo~\ref{cap:projetos} para criar o primeiro projeto.

Cerca de seis segundos após abrir, a AURORA verifica silenciosamente se há uma versão mais nova; se houver, você verá a janela de atualização (Capítulo~\ref{cap:atualizacoes}). Não é preciso fazer nada para manter o SAPHO em dia além de aceitar as atualizações quando oferecidas.

\section{Desinstalação}

Use \menu{Configurações do Windows $\to$ Aplicativos $\to$ Aurora-IDE $\to$ Desinstalar}. Os projetos criados por você não são apagados --- eles pertencem às pastas que você escolheu.

% !TEX root = ../main.tex
\chapter{Tour pela interface}
\label{cap:tour-interface}

A janela da AURORA não tem a moldura padrão do Windows nem uma barra de menus \menu{Arquivo/Editar/Exibir}: tudo vive numa \textbf{barra de título personalizada} (que concentra a toolbar), no \textbf{painel lateral} da árvore de arquivos, no \textbf{editor}, nos \textbf{terminais} e na \textbf{barra de status}. Este capítulo apresenta cada região; os detalhes de uso ficam nos capítulos dedicados.

\figurapendente{Janela principal da AURORA com um projeto aberto}{Visão geral anotável: toolbar no topo, árvore à esquerda, editor ao centro, terminais embaixo, barra de status no rodapé.}

\section{A barra de título e a toolbar}
\label{sec:toolbar}

A barra superior tem três zonas:

\textbf{Zona esquerda} --- identidade e projeto:
\begin{itemize}
  \item o ícone do SAPHO;
  \item \botao{Alternar barra lateral} --- mostra/esconde o painel da árvore de arquivos;
  \item \botao{Novo Projeto} e \botao{Abrir Projeto}.
\end{itemize}

\textbf{Zona central} --- a cadeia do processador, na ordem em que você a usa:
\begin{itemize}
  \item \botao{Hub de Processadores} --- cria processadores (Capítulo~\ref{cap:criando-processadores});
  \item \botao{Compilar \cmm{}} (glifo \enquote{C±}) --- compila o fonte \cmm{} até Verilog;
  \item \botao{Configurações de simulação do processador} (engrenagem) --- clock, número de ciclos e exibição de arrays;
  \item \botao{Sintetizar Verilog} (pena) --- verifica os arquivos Verilog com o Icarus Verilog;
  \item \botao{PRISM} --- abre o visualizador de RTL (Capítulo~\ref{cap:prism});
  \item \textbf{seletor de simulador} --- Icarus Verilog ou Verilator (Capítulo~\ref{cap:simulacao});
  \item \textbf{seletor de visualizador} --- GTKWave ou Surfer (Capítulo~\ref{cap:ondas});
  \item \botao{Analisar Verilog} (onda quadrada) --- roda a simulação e abre as formas de onda;
  \item \botao{Fast Sim} (\enquote{»}) --- simulação rápida sem ondas (apenas com Verilator);
  \item \botao{Cancelar} --- interrompe todos os processos de compilação/simulação em andamento;
  \item \botao{Configuração de ondas} --- escolhe quais sinais entram no arquivo de ondas;
  \item \textbf{seletor de \file{.gtkw}} --- escolhe o layout de ondas ativo;
  \item \botao{Teste do processador sintetizado} (chip) --- teste de E/S do processador via Verilator.
\end{itemize}

\textbf{Zona direita} --- apoio e janela:
\begin{itemize}
  \item \botao{Controle de Versão} (Dagr, Capítulo~\ref{cap:git});
  \item \botao{Configurações do Aurora} (Capítulo~\ref{cap:configuracoes});
  \item \botao{Assistente IA} (estrela --- Aurora Intelligence, Capítulo~\ref{cap:aurora-intelligence});
  \item os controles de janela: minimizar, maximizar/restaurar, fechar. Duplo clique na área vazia da barra também maximiza/restaura.
\end{itemize}

\figurapendente{Toolbar completa em close}{Print da barra de título inteira, com um projeto aberto e os botões habilitados. Idealmente uma segunda versão com os tooltips de um dos botões visível.}

\section{O painel da árvore de arquivos}

À esquerda fica o painel \textbf{ÁRVORE DE ARQUIVOS}, com o nome do projeto aberto ao lado do rótulo. Ele tem \textbf{três visões}, alternadas por um único botão no cabeçalho do painel:

\begin{enumerate}
  \item \textbf{Arquivos} --- a visão padrão, agrupada por processador, com os arquivos Verilog classificados automaticamente entre sintetizáveis e \emph{testbench};
  \item \textbf{Hierarquia} --- a árvore de instâncias de módulos do projeto, extraída pelo Yosys após uma síntese;
  \item \textbf{Pastas} --- um explorador de diretórios convencional, no estilo do VS Code, com criar/renomear/mover/excluir.
\end{enumerate}

O cabeçalho do painel traz ainda: \botao{Novo arquivo}, \botao{Atualizar}, \botao{Buscar nos arquivos} (\tecla{Ctrl+Shift+F}), \botao{Abrir no explorador do sistema}, \botao{Backup do projeto} (gera um \file{.zip}), \botao{Recolher tudo} e \botao{Fechar projeto}. O uso completo das árvores está no Capítulo~\ref{cap:projetos}.

\section{O editor}

O centro da janela é o editor de código, baseado no \tool{Monaco} (o mesmo motor do VS Code): abas múltiplas, divisão em até três painéis, realce de sintaxe para \cmm{}, assembly, Verilog e SystemVerilog, diagnósticos em tempo real e formatação automática. Tudo no Capítulo~\ref{cap:editor}.

\section{Os terminais}

A área inferior reúne \textbf{seis terminais} em abas: \textbf{TCMM}, \textbf{TASM}, \textbf{TVERI}, \textbf{TWAVE} e \textbf{THTEST} são consoles de saída das etapas de compilação e simulação (cada etapa escreve no seu), e o \textbf{TCMD} é um terminal PowerShell interativo de verdade. Filtros por tipo de mensagem, exportação de log e números de linha clicáveis completam o conjunto (Capítulo~\ref{cap:terminais}).

\section{A barra de status}

No rodapé:
\begin{itemize}
  \item à esquerda, o indicador \textbf{Pronto/Não Pronto} (o ponto verde/vermelho indica se o projeto está em condições de compilar) e o \textbf{processador ativo} (deduzido do arquivo \file{.cmm} em foco no editor);
  \item ao centro, o botão \botao{Iniciar Compilação} (ícone de raio);
  \item à direita, avisos de estrutura (\enquote{Sem top-level}, \enquote{Sem testbench}), o motor de simulação selecionado, a posição do cursor (\code{Ln}, \code{Col}) e o indicador do GitHub/controle de versão.
\end{itemize}

\section{Paleta de comandos, notificações e tooltips}

\begin{itemize}
  \item \textbf{Paleta de comandos} (\tecla{Ctrl+Shift+P}): busca difusa sobre os comandos da IDE, organizados em grupos (Compile, Project, View, Tools). Útil quando você sabe o nome do que quer fazer, mas não onde fica o botão.
  \item \textbf{Notificações}: eventos como criação de arquivos, troca de simulador e avisos de projeto aparecem como \emph{toasts} no canto da janela.
  \item \textbf{Tooltips}: todo botão tem uma dica de ferramenta rica ao passar o mouse; elas podem ser desligadas nas Configurações.
\end{itemize}

\section{A tela de boas-vindas}

Sem projeto aberto, a AURORA mostra a tela de boas-vindas: a marca SAPHO com a \emph{tagline} \enquote{Processador de Arquitetura Escalável para Otimização de Hardware}, a seção \textbf{Início} (\botao{Novo Projeto}, \botao{Abrir Projeto}) e a seção \textbf{Recentes} com os últimos dez projetos. O fundo é uma animação de aurora boreal.


\part{Projetos e processadores}
% !TEX root = ../main.tex
\chapter{Projetos}
\label{cap:projetos}

Tudo no SAPHO acontece dentro de um \textbf{projeto}: uma pasta no disco governada por um arquivo \file{.spf} (\emph{SAPHO Project File}). O projeto reúne um ou mais processadores, os arquivos Verilog de integração (\emph{top-level} e \emph{testbench}) e os metadados que a AURORA usa para orquestrar compilação e simulação.

\section{Criando um projeto}
\label{sec:criando-projeto}

Clique em \botao{Novo Projeto} (na toolbar ou na tela de boas-vindas). O modal \textbf{Criar Novo Projeto} pede:

\begin{longtable}{@{}lp{0.64\linewidth}@{}}
\toprule
\textbf{Campo} & \textbf{Regras} \\
\midrule
\endhead
\textbf{Nome do Projeto} & Apenas letras, números, sublinhado (\code{\_}) e hífen (\code{-}). Sem espaços, acentos ou símbolos. \\
\textbf{Local do Projeto} & Somente leitura; preencha com o botão \botao{Procurar}, que abre o seletor de pastas do Windows. A AURORA lembra o último local usado. \\
\bottomrule
\end{longtable}

Clique em \botao{Gerar Projeto}. Se algum campo violar as regras, o diálogo \textbf{Entrada Inválida} explica o problema (\enquote{Nome do Projeto e Local não podem conter espaços ou símbolos especiais\ldots}).

\figurapendente{Modal Criar Novo Projeto preenchido}{Com o nome \code{MeuFiltro} e um local qualquer, antes de clicar em Gerar Projeto.}

No disco, é criada a estrutura \file{<Local>/<Nome>/<Nome>.spf}. Um projeto recém-criado ainda não tem processadores --- a árvore abre praticamente vazia, e o próximo passo natural é o \botao{Hub de Processadores} (Capítulo~\ref{cap:criando-processadores}).

Para o exemplo condutor deste manual, crie o projeto \textbf{MeuFiltro}.

\section{O arquivo \texttt{.spf}}
\label{sec:spf}

O \file{.spf} é um arquivo JSON legível, com dois blocos:

\begin{itemize}
  \item \code{metadata} --- nome do projeto, datas de criação/modificação/última abertura, nome do computador e versão da AURORA que o criou;
  \item \code{structure} --- a lista de \code{processors}, as pastas do projeto, o \code{topLevelFile}, o \code{testbenchFile} e as listas de arquivos Verilog sintetizáveis e de \emph{testbench}.
\end{itemize}

Você raramente precisará editá-lo à mão: a AURORA é a única escritora do arquivo e o mantém em dia a cada ação na interface. Os caminhos são gravados relativos à pasta do projeto, então o projeto inteiro pode ser copiado, movido de máquina ou versionado no Git sem quebrar --- ao abrir, a AURORA realinha os caminhos ao novo local.

\section{Abrindo um projeto}

Três caminhos:
\begin{itemize}
  \item \botao{Abrir Projeto} na toolbar ou na tela de boas-vindas, selecionando o \file{.spf};
  \item dois cliques no \file{.spf} pelo Explorer do Windows (a extensão é associada na instalação);
  \item a lista de \textbf{Recentes} na tela de boas-vindas (até dez projetos, também disponíveis na \emph{jumplist} do ícone da AURORA na barra de tarefas do Windows).
\end{itemize}

Ao abrir, a AURORA carrega a árvore, restaura a lista de processadores e registra a data de abertura. Se o \file{.spf} tiver sumido do caminho registrado, o diálogo \textbf{Projeto Não Encontrado} avisa e o item sai da lista de recentes.

\section{Arquivos do projeto: quem é quem}
\label{sec:arquivos-do-projeto}

Um projeto maduro tem esta cara (exemplo com o nosso \code{MeuFiltro}):

\begin{lstlisting}
MeuFiltro/
|-- MeuFiltro.spf              (o arquivo de projeto)
|-- media_movel/               (um processador)
|   |-- Software/media_movel.cmm
|   |-- Hardware/              (gerados: media_movel.v, *.mif)
|   `-- Simulation/            (gerado: media_movel_tb.v)
|-- TopLevel/                  (Verilog de integracao, opcional)
`-- ...
\end{lstlisting}

Dois papéis especiais são atribuídos por você, pelo menu de contexto da árvore:

\begin{itemize}
  \item \textbf{Top-level} --- o módulo Verilog que integra o(s) processador(es); é o topo da síntese e o alvo do PRISM. Menu de contexto: \menu{Definir como Top Level} / \menu{Remover Top Level}.
  \item \textbf{Testbench} --- o banco de testes da simulação (um \file{_tb.v} Verilog ou um \file{.py} cocotb). Menu de contexto: \menu{Marcar como Testbench} / \menu{Desmarcar como Testbench}.
\end{itemize}

A barra de status lembra você do que falta: \enquote{Sem top-level}, \enquote{Sem testbench}. Para um projeto de um único processador, o \emph{testbench} gerado automaticamente pela compilação (\file{<proc>_tb.v}) costuma bastar --- os papéis manuais importam quando você escreve integração própria.

\section{As três visões da árvore}
\label{sec:tres-arvores}

O botão de alternância no cabeçalho do painel circula entre as visões:

\subsection{Visão Arquivos}
A visão padrão: um \emph{picker} agrupado por processador, mostrando os fontes \cmm{}, os Verilog classificados automaticamente (sintetizável vs.\ \emph{testbench}) e os artefatos relevantes. É a visão de trabalho do dia a dia.

\subsection{Visão Hierarquia}
Depois de uma síntese (por exemplo, ao abrir o PRISM), esta visão mostra a \textbf{árvore de instâncias de módulos} do projeto, como o Yosys a enxerga: o top-level no topo, cada instância aninhada abaixo. Só aparece quando há dados de hierarquia.

\subsection{Visão Pastas}
Um explorador de arquivos completo, no estilo do VS Code, com CRUD integral pelo menu de contexto:

\begin{itemize}
  \item \menu{Novo Arquivo\ldots} / \menu{Nova Pasta\ldots} / \menu{Novo arquivo Verilog (.v)} / \menu{Novo testbench cocotb (.py)} / \menu{Novo .gitignore};
  \item \menu{Recortar}, \menu{Copiar}, \menu{Colar}, \menu{Renomear\ldots}, \menu{Excluir} (com escolha entre \menu{Mover para a Lixeira} --- restaurável --- e \menu{Excluir Permanentemente});
  \item \menu{Copiar Caminho} / \menu{Copiar Caminho Relativo};
  \item \menu{Abrir no Terminal Integrado} (abre o TCMD já na pasta) e \menu{Revelar no Explorador de Arquivos};
  \item \menu{Atualizar} e \menu{Recolher Tudo}.
\end{itemize}

Nomes de arquivo são validados como o Windows exige: sem \texttt{< > : " | ? *}, sem nomes reservados do sistema, sem terminar em ponto ou espaço. Em conflitos de colagem, a AURORA oferece \menu{Substituir} ou \menu{Manter Ambos}.

\figurapendente{As três visões da árvore, lado a lado}{Três prints da mesma região: visão Arquivos, visão Hierarquia (após abrir o PRISM) e visão Pastas com o menu de contexto aberto.}

\begin{nota}
\textbf{Arrastar e soltar:} a árvore aceita a importação de arquivos \file{.v}, \file{.sv}, \file{.vh} e \emph{testbenches} cocotb \file{.py} por arrastar-e-soltar. Outros tipos são recusados com uma dica (arquivos \file{.gtkw}, por exemplo, entram pelo seletor próprio da toolbar).
\end{nota}

\section{Manutenção do projeto}

\begin{itemize}
  \item \textbf{Backup}: o botão \botao{Backup do projeto} no cabeçalho da árvore gera um \file{.zip} do projeto inteiro, sob demanda.
  \item \textbf{Arquivos faltantes}: se arquivos registrados no \file{.spf} sumirem do disco, um aviso lista quantos faltam e oferece dispensá-los (removendo as referências).
  \item \textbf{Fechar projeto}: o botão vermelho no cabeçalho da árvore, com confirmação (\enquote{Tem certeza que deseja fechar o projeto atual?}).
  \item \textbf{Renomear projeto}: disponível também via Aurora Intelligence; a operação reescreve o \file{.spf} e a pasta.
\end{itemize}

% !TEX root = ../main.tex
\chapter{O processador SAPHO}
\label{cap:processador-sapho}

Antes de criar processadores na IDE, vale entender o que exatamente está sendo criado. Este capítulo descreve a arquitetura do \emph{soft-processor} SAPHO do ponto de vista de quem o usa: o que os parâmetros significam, como o processador executa o seu programa e quais são as suas restrições. Não é preciso decorar nada aqui para usar a plataforma --- mas quem entende a máquina escreve programas melhores para ela.

\section{Visão geral da arquitetura}

O SAPHO é um processador de \textbf{acumulador} com \textbf{arquitetura Harvard}:

\begin{itemize}
  \item \textbf{Acumulador}: em vez de um banco de registradores, existe um único registrador central, o acumulador, que recebe o resultado de toda operação da ULA (unidade lógico-aritmética). O padrão dominante é: um operando vem da memória de dados, o outro é o próprio acumulador, e o resultado volta para o acumulador. Uma \textbf{pilha de dados} fornece o segundo operando quando a expressão exige.
  \item \textbf{Harvard}: as memórias de programa e de dados são fisicamente separadas, cada uma com sua largura e seu barramento. O programa não pode se sobrescrever, e dados e instruções são acessados em paralelo.
\end{itemize}

Três traços definem o estilo do processador:

\begin{enumerate}
  \item \textbf{Totalmente parametrizável} --- toda largura é um parâmetro Verilog: a palavra de dados (\code{NUBITS}), a mantissa (\code{NBMANT}) e o expoente (\code{NBEXPO}) do ponto flutuante, o operando das instruções, as profundidades das memórias e das pilhas.
  \item \textbf{Pague pelo que usar} --- cada instrução do conjunto é um parâmetro booleano; o montador do YANC liga apenas as instruções que o seu programa realmente usa, e o hardware das demais é eliminado na geração do circuito.
  \item \textbf{Ciclo curto} --- o processador opera com um \emph{pipeline} raso de três estágios (busca, decodificação, execução), sem \emph{forwarding}: a ULA é combinacional e o resultado se acumula a cada ciclo, de modo que nem saltos condicionais quebram o fluxo. O comportamento temporal é simples de prever, o que é uma virtude em instrumentação.
\end{enumerate}

\section{Os parâmetros fundamentais}
\label{sec:parametros-fundamentais}

Ao criar um processador (Capítulo~\ref{cap:criando-processadores}) você escolhe, direta ou indiretamente, os valores abaixo. Eles aparecem como diretivas no topo do arquivo \cmm{} e como \code{parameter} no Verilog gerado.

\begin{longtable}{@{}llp{0.60\linewidth}@{}}
\toprule
\textbf{Parâmetro} & \textbf{Diretiva \cmm{}} & \textbf{Significado} \\
\midrule
\endhead
\code{NUBITS} & \code{\#NUBITS} & Largura da palavra de dados, em bits. Todo inteiro do programa tem esse tamanho, em complemento de dois. \\
\code{NBMANT} & \code{\#NBMANT} & Largura da mantissa do ponto flutuante. \\
\code{NBEXPO} & \code{\#NBEXPO} & Largura do expoente do ponto flutuante. \\
--- & \code{\#NDSTAC} & Profundidade da pilha de dados. \\
--- & \code{\#SDEPTH} & Profundidade da pilha de sub-rotinas (chamadas de função aninhadas). \\
--- & \code{\#NUIOIN} & Número de portas de entrada. \\
--- & \code{\#NUIOOU} & Número de portas de saída. \\
--- & \code{\#NUGAIN} & Constante usada pela função \code{norm()} (normalização por divisão fixa). \\
--- & \code{\#FFTSIZ} & Tamanho da FFT ($2^{\texttt{FFTSIZ}}$ pontos), quando o programa usa os endereçamentos com reversão de bits. \\
\bottomrule
\end{longtable}

\begin{aviso}
\textbf{Restrição estrutural do ponto flutuante:} a palavra precisa comportar exatamente o formato, isto é, \texttt{NUBITS} = \texttt{NBMANT} + \texttt{NBEXPO} + 1 (o bit extra é o sinal). A AURORA valida essa igualdade no formulário de criação de processador e recusa combinações inválidas.
\end{aviso}

\section{O ponto flutuante do SAPHO}
\label{sec:float-sapho}

O SAPHO tem hardware de ponto flutuante \textbf{próprio, que não segue o padrão IEEE 754}. O formato é: 1 bit de sinal, \code{NBEXPO} bits de expoente em complemento de dois e \code{NBMANT} bits de mantissa com o bit líder explícito. Isso permite escolher precisões arbitrárias --- uma mantissa de 10 bits com expoente de 5, por exemplo --- e economizar hardware exatamente onde a aplicação permite.

Na prática, para o usuário:
\begin{itemize}
  \item Dentro do programa \cmm{}, \code{float} \enquote{simplesmente funciona}: literais, aritmética e as funções da biblioteca operam nesse formato interno.
  \item A conversão entre o formato IEEE do seu computador e o formato do SAPHO é feita em software pela toolchain (na carga das memórias e na leitura dos resultados), de forma transparente.
  \item A precisão que você obtém é exatamente a que você configurou: com \code{\#NBMANT 10}, espere cerca de 3 dígitos decimais significativos.
\end{itemize}

Números \textbf{complexos} (tipo \code{comp} da linguagem) não têm hardware dedicado: cada complexo é um par de \emph{floats} (parte real e imaginária) manipulado componente a componente pelo código gerado.

\section{Como o programa vira circuito}

O processador gerado é composto por poucos módulos Verilog, todos criados a partir de moldes parametrizáveis da pasta de HDL do YANC:

\begin{longtable}{@{}lp{0.68\linewidth}@{}}
\toprule
\textbf{Módulo} & \textbf{Função} \\
\midrule
\endhead
\code{processor} & O topo: instancia o núcleo e as duas memórias. \\
\code{mem\_instr} & Memória de programa (somente leitura), carregada com a imagem \file{<proc>_inst.mif}. \\
\code{mem\_data} & Memória de dados, carregada com a imagem \file{<proc>_data.mif}; uma leitura e uma escrita por ciclo. \\
\code{core} & O caminho de dados: acumulador, pilhas, busca de instruções, contador de programa. \\
\code{instr\_dec} & Decodificador de instruções. \\
\code{ula} & A unidade lógico-aritmética: um submódulo por operação, e só os usados sobrevivem. \\
\code{addr\_dec} & Decodificador de endereços para múltiplas portas de E/S. \\
\code{myFIFO} & FIFO genérica, disponível para acoplar o processador ao mundo externo sem depender de IP de fabricante. \\
\bottomrule
\end{longtable}

O conjunto de instruções usa opcode de 7 bits, com famílias para carga e armazenamento (direto, indireto e com reversão de bits para FFT), pilha, entrada e saída, controle de fluxo (salto incondicional, salto se zero, chamada e retorno de sub-rotina) e cálculo (aritmética inteira e em ponto flutuante, lógica, comparações, deslocamentos e conversões). Você não escreve essas instruções à mão --- o compilador \cmm{} as gera ---, mas as verá nas formas de onda: a AURORA exibe, em \emph{lockstep} com o clock, a instrução assembly e a linha \cmm{} correspondente a cada ciclo (Capítulo~\ref{cap:ondas}).

\section{Entrada, saída e o mundo externo}

O processador conversa com o exterior por portas numeradas de entrada e de saída, criadas conforme \code{\#NUIOIN} e \code{\#NUIOOU}. No programa, elas são acessadas pelas funções \code{in()}, \code{out()}, \code{fin()} e \code{fout()} (Capítulo~\ref{cap:cmm-io-stdlib}). No hardware, cada porta vira sinais de dados e de controle no topo do módulo \code{processor} --- é por eles que o seu projeto de FPGA conecta o processador a ADCs, DACs, FIFOs e outros blocos.

Dois recursos opcionais completam a interface: \textbf{interrupção} (o processador desvia para um endereço configurado quando o pino de interrupção pulsa) e o marcador \code{\#TOAQUI}, que gera um pino \code{cheguei} pulsado quando o programa passa por um ponto marcado --- útil para sincronizar hardware externo com fases do programa.

\section{Um processador por tarefa}

Como cada processador é gerado sob medida, o padrão de projeto no SAPHO é diferente do mundo dos processadores fixos: em vez de um processador grande que faz tudo, criam-se \textbf{vários processadores pequenos}, um por tarefa, conectados entre si por um \emph{top-level} Verilog (por exemplo, um processador que detecta cruzamentos de zero alimentando outro que calcula uma distância). A AURORA suporta esse fluxo diretamente: um projeto pode conter quantos processadores você quiser, cada um com sua pasta, seu \file{.cmm} e seus parâmetros, todos integrados por um arquivo \emph{top-level} e um \emph{testbench} comuns (Capítulo~\ref{cap:projetos}).

% !TEX root = ../main.tex
\chapter{Criando processadores: o Hub de Processadores}
\label{cap:criando-processadores}

Com o projeto aberto, o botão \botao{Hub de Processadores} na toolbar abre o formulário que dá vida a um novo processador SAPHO. É aqui que as escolhas de arquitetura do Capítulo~\ref{cap:processador-sapho} viram campos concretos.

\figurapendente{Modal Criar Hub de Processadores}{O formulário completo, preenchido para o processador \code{media\_movel} do exemplo condutor (Total de Bits 16, Mantissa 10, Expoente 5).}

\section{Os campos do formulário}

\begin{longtable}{@{}p{0.30\linewidth}llp{0.34\linewidth}@{}}
\toprule
\textbf{Campo} & \textbf{Diretiva} & \textbf{Padrão} & \textbf{Validação} \\
\midrule
\endhead
Nome do Processador & \code{\#PRNAME} & \code{procTest\_00} & Letras, números, \code{\_} e \code{-}; sem espaços ou acentos. \\
\multicolumn{4}{@{}l}{\emph{Seção Formato Numérico}} \\
Total de Bits & \code{\#NUBITS} & 23 & Inteiro $>0$ e igual a Mantissa + Expoente + 1. \\
Ganho & \code{\#NUGAIN} & 128 & Inteiro $>0$ e \textbf{potência de 2}. \\
Bits da Mantissa & \code{\#NBMANT} & 16 & Inteiro $>0$. \\
Bits do Expoente & \code{\#NBEXPO} & 6 & Inteiro $>0$. \\
\multicolumn{4}{@{}l}{\emph{Seção Pilhas e Portas}} \\
Tamanho da Pilha de Instruções & \code{\#SDEPTH} & 5 & Inteiro $>0$. \\
Tamanho da Pilha de Dados & \code{\#NDSTAC} & 5 & Inteiro $>0$. \\
Portas de Entrada & \code{\#NUIOIN} & 1 & Inteiro $>0$. \\
Portas de Saída & \code{\#NUIOOU} & 1 & Inteiro $>0$. \\
\bottomrule
\end{longtable}

O formulário valida enquanto você digita: um campo inválido ganha borda vermelha e um tooltip explicando a regra (\enquote{Total de Bits deve ser igual a Bits da Mantissa + Bits do Expoente + 1}, \enquote{Ganho deve ser potência de 2}). O botão \botao{Gerar Processador} só habilita quando tudo está válido. Note que o padrão de fábrica já é consistente: $23 = 16 + 6 + 1$.

Como escolher os valores:
\begin{itemize}
  \item \textbf{Total de Bits} dimensiona os inteiros e o consumo de hardware; 16 bits bastam para muitos sinais de instrumentação, 23 (padrão) dá folga com \emph{float} de boa precisão.
  \item \textbf{Mantissa/Expoente} definem a precisão e a faixa do ponto flutuante (Seção~\ref{sec:float-sapho}); se o seu programa só usa inteiros, os valores têm pouco efeito prático, mas a igualdade estrutural continua obrigatória.
  \item \textbf{Pilha de Instruções} limita a profundidade de chamadas de função aninhadas; \textbf{Pilha de Dados} limita a complexidade das expressões. Os padrões atendem programas típicos.
  \item \textbf{Portas} de entrada/saída: uma por fluxo de dados que o processador troca com o mundo (Capítulo~\ref{cap:cmm-io-stdlib}).
  \item \textbf{Ganho}: o divisor fixo da função \code{norm()}; relevante apenas se você a usa.
\end{itemize}

\section{O que é gerado}

Ao clicar em \botao{Gerar Processador}, a AURORA cria dentro do projeto:

\begin{lstlisting}
<projeto>/<NomeDoProcessador>/
|-- Software/<NomeDoProcessador>.cmm
|-- Hardware/        (vazia; recebera o Verilog e as memorias)
`-- Simulation/      (vazia; recebera o testbench)
\end{lstlisting}

O \file{.cmm} inicial vem com o cabeçalho de diretivas preenchido com os valores do formulário e um \code{main()} vazio à sua espera:

\begin{lstlisting}[style=cmm]
#PRNAME media_movel
#NUBITS 16
#NDSTAC 5
#SDEPTH 5
#NUIOIN 1
#NUIOOU 1
#NBMANT 10
#NBEXPO 5
#NUGAIN 128

void main()
{
    // Øk. Você criou um processador em C±, mas e agora?
}
\end{lstlisting}

E agora, o Capítulo~\ref{cap:cmm-fundamentos}: escrever o programa. O processador também é registrado no \file{.spf} (nomes duplicados são bloqueados, sem distinção de maiúsculas).

\section{Configurações de simulação do processador}
\label{sec:proc-config}

Ao lado do botão de compilar há a engrenagem \botao{Configurações de simulação do processador}, que abre um popover com três ajustes \textbf{por processador}:

\begin{longtable}{@{}lp{0.52\linewidth}l@{}}
\toprule
\textbf{Campo} & \textbf{Efeito} & \textbf{Padrão} \\
\midrule
\endhead
\textbf{Clock frequency (MHz)} & Frequência do clock do \emph{testbench} gerado. & 100 \\
\textbf{Number of clocks} & Quantos ciclos a simulação executa antes de encerrar. & 2000 \\
\textbf{Show arrays in GTKWave} & Emite cada elemento de array como um sinal visível nas ondas (equivale à flag \code{--array} do compilador). & desligado \\
\bottomrule
\end{longtable}

O popover também exibe o \textbf{tempo estimado de simulação} (ciclos divididos pelo clock, em $\mu$s). Se a sua simulação termina antes de o programa produzir os resultados, aumente o número de ciclos aqui.

\section{Renomear e excluir processadores}

Pelo menu de contexto do processador na árvore. Renomear é uma operação completa: a pasta é movida, os artefatos (\file{.cmm}, \file{.asm}, \file{Hardware/<nome>.v}, \file{Simulation/<nome>_tb.v}) são renomeados e a diretiva \code{\#PRNAME} dentro do fonte é corrigida --- tudo numa ação. As mesmas regras de nome do formulário se aplicam. Excluir remove a pasta do processador e o registro no \file{.spf}.

\begin{aviso}
O nome do processador e a diretiva \code{\#PRNAME} precisam concordar --- é assim que a toolchain conecta o fonte aos artefatos. Renomeie sempre pela AURORA, que mantém os dois lados em sincronia.
\end{aviso}


\part{A linguagem \cmm{}}
% !TEX root = ../main.tex
\chapter{\cmm{}: fundamentos}
\label{cap:cmm-fundamentos}

A \cmm{} (lê-se \enquote{C mais-menos}; nos arquivos, extensão \file{.cmm}) é a linguagem principal do SAPHO: um dialeto de C criado pelo NIPS-CERN, enxuto o bastante para virar hardware eficiente e expressivo o bastante para algoritmos de processamento de sinais --- incluindo números complexos como tipo nativo e álgebra linear em notação de Dirac. Este capítulo apresenta a estrutura de um programa; os capítulos seguintes cobrem tipos, operadores, E/S, biblioteca e recursos avançados.

\begin{nota}
Tudo o que está nesta Parte foi extraído da gramática oficial do compilador (\file{CMMComp.y} e \file{CMMComp.l} do repositório \repo{yanc}, toolchain 5.3). Quando este manual diz que algo \enquote{não existe} na linguagem, é porque não existe na gramática --- o compilador rejeitará o código.
\end{nota}

\section{A estrutura de um programa}

Um programa \cmm{} é uma sequência de três tipos de elemento, em qualquer ordem: \textbf{diretivas} (linhas iniciadas por \code{\#}), \textbf{declarações de variáveis globais} e \textbf{funções}. O ponto de entrada é obrigatório: uma função \code{void main()}. Sem ela, o compilador encerra com erro.

A forma canônica de um programa SAPHO é: diretivas no topo (configurando o processador), depois o \code{main()} com um laço infinito \code{while (1)} --- afinal, o programa modela um circuito que roda continuamente, processando amostras para sempre:

\begin{lstlisting}[style=cmm]
#PRNAME media_movel
#NUBITS 16
#NDSTAC 4
#SDEPTH 2
#NUIOIN 1
#NUIOOU 1
#NBMANT 10
#NBEXPO 5

void main()
{
    int x[4];    // historico das ultimas 4 amostras
    int soma;

    while (1)
    {
        x[3] = x[2];          // desloca o historico
        x[2] = x[1];
        x[1] = x[0];
        x[0] = in(0);         // le nova amostra da porta 0

        soma = x[0] + x[1] + x[2] + x[3];
        out(0, soma >> 2);    // media = soma/4, sem usar divisor
    }
}
\end{lstlisting}

Este é o \code{media_movel}, o processador que acompanha o manual inteiro: um filtro de média móvel de 4 amostras. Repare em duas escolhas típicas de quem escreve para o SAPHO: o histórico é um array deslocado à mão, e a divisão por 4 é um deslocamento de bits (\code{>> 2}) --- assim o processador gerado não precisa instanciar um divisor, que é um dos blocos mais caros da ULA.

\section{Diretivas: configurando o processador no próprio código}
\label{sec:diretivas}

As diretivas são linhas que começam com \code{\#}, uma por linha, \textbf{sem ponto e vírgula}. Elas definem o nome e a arquitetura do processador que será gerado. Quando você cria um processador pela AURORA (Capítulo~\ref{cap:criando-processadores}), o formulário preenche essas diretivas por você no \file{.cmm} inicial; depois disso, o arquivo é a fonte da verdade --- editar a diretiva muda o processador na próxima compilação.

\begin{longtable}{@{}llp{0.52\linewidth}@{}}
\toprule
\textbf{Diretiva} & \textbf{Padrão} & \textbf{Significado} \\
\midrule
\endhead
\code{\#PRNAME <nome>} & --- & Nome do processador (deve casar com o nome usado no projeto). \\
\code{\#NUBITS <n>} & 23 & Largura da palavra em bits (inteiros em complemento de dois). \\
\code{\#NBMANT <n>} & 16 & Bits de mantissa do ponto flutuante. \\
\code{\#NBEXPO <n>} & 6 & Bits de expoente do ponto flutuante. \\
\code{\#NDSTAC <n>} & 10 & Profundidade da pilha de dados. \\
\code{\#SDEPTH <n>} & 10 & Profundidade da pilha de sub-rotinas (limita chamadas aninhadas). \\
\code{\#NUIOIN <n>} & 1 & Número de portas de entrada. \\
\code{\#NUIOOU <n>} & 1 & Número de portas de saída. \\
\code{\#NUGAIN <n>} & 64 & Divisor fixo usado pela função \code{norm()}. \\
\code{\#FFTSIZ <n>} & 8 & Tamanho da FFT ($2^n$ pontos) para o endereçamento bit-reverso. \\
\bottomrule
\end{longtable}

\begin{aviso}
Lembre da restrição estrutural: \texttt{NUBITS} = \texttt{NBMANT} + \texttt{NBEXPO} + 1. No exemplo acima, $16 = 10 + 5 + 1$. A AURORA valida isso ao criar o processador, mas se você editar as diretivas à mão, é sua responsabilidade manter a igualdade.
\end{aviso}

Existem ainda duas diretivas \emph{comportamentais}, que aparecem no meio do código (não no topo) e são explicadas no Capítulo~\ref{cap:cmm-avancado}: \code{\#PRACA} (ponto de retorno da interrupção) e \code{\#TOAQUI} (marcador refletido no pino \code{cheguei}).

\section{Constantes com \texttt{\#define}}

A \cmm{} tem um pré-processador mínimo embutido: \code{\#define NOME corpo} define uma constante simbólica, expandida onde o nome aparecer.

\begin{lstlisting}[style=cmm]
#define SCALE  3
#define OFFSET 7

void main()
{
    int x;
    while (1)
    {
        x = in(0);
        x = x + SCALE;
        out(0, x);
        x = x + OFFSET;
        out(0, x);
    }
}
\end{lstlisting}

Três limitações importantes em relação ao C tradicional:
\begin{itemize}
  \item só existem \emph{defines de objeto} --- macros com argumentos, como \code{\#define DOBRO(x) ((x)*2)}, não são suportadas;
  \item não há \code{\#include} nem \code{\#ifdef} no lado \cmm{} (o fluxo C++, sim, tem pré-processador completo --- Capítulo~\ref{cap:cmm-avancado});
  \item o limite é de 256 defines por programa.
\end{itemize}

O editor da AURORA reconhece os seus \code{\#define} e os colore como constantes conforme você digita.

\section{Comentários}

Como em C: \code{//} comenta até o fim da linha, \code{/* ... */} comenta blocos.

\section{Onde o arquivo vive}

Dentro de um projeto AURORA, o código de cada processador fica em \file{<projeto>/<processador>/Software/<processador>.cmm}. A compilação gera o assembly em \file{Software/}, o Verilog e as imagens de memória em \file{Hardware/} e o \emph{testbench} em \file{Simulation/} (Capítulo~\ref{cap:compilacao}).

% !TEX root = ../main.tex
\chapter{\cmm{}: tipos, operadores e controle}
\label{cap:cmm-tipos-operadores}

\section{Tipos de dados}

A \cmm{} tem exatamente quatro palavras-chave de tipo. O \code{void} marca a ausência de tipo, no retorno de funções que não retornam nada. O \code{int} é o inteiro de \code{\#NUBITS} bits, em complemento de dois. O \code{float} é o ponto flutuante no formato próprio do SAPHO, com mantissa de \code{\#NBMANT} bits e expoente de \code{\#NBEXPO} bits (Seção~\ref{sec:float-sapho}). E o \code{comp} é o número complexo, um par de \code{float} com as partes real e imaginária, tratado como tipo de primeira classe da linguagem.

Não existe um tipo de ponto fixo separado, papel que cabe ao \code{float} configurável, e tampouco existem \code{char}, \code{double}, \code{struct}, \code{union}, \code{enum} ou \textit{strings}; um literal de texto só aparece na inicialização de \textit{arrays} por arquivo, adiante.

\subsection{Literais complexos e o \texttt{i} reservado}

Um literal complexo escreve-se \code{a+bi}, como em matemática:

\begin{lstlisting}[style=cmm]
comp z;
z = 3.0 + 4.0i;
\end{lstlisting}

Por causa disso, o identificador \code{i} é reservado para a unidade imaginária, e usá-lo como nome de variável é erro de compilação. Adote \code{k}, \code{n} ou \code{idx} para contadores.

\subsection{Conversões entre tipos}

Misturar \code{int}, \code{float} e \code{comp} numa expressão ou atribuição é permitido, mas o compilador emite avisos de conversão implícita, por exemplo ao guardar um \code{float} num \code{int} com perda. Usar \code{float} ou \code{comp} como condição de \code{if} ou \code{while}, ou como índice de \textit{array}, também gera aviso. Trate esses avisos como cheiro de projeto: conversões explícitas e índices inteiros deixam o hardware gerado mais barato e o comportamento mais óbvio.

\section{Operadores}

A Tabela~\ref{tab:operadores} reúne os operadores da linguagem, cuja precedência segue a do C.

\begin{table}[ht]
\centering
\caption{Operadores da \cmm{}.}
\label{tab:operadores}
\small
\begin{tabular}{lp{0.56\linewidth}}
\toprule
Grupo & Operadores \\
\midrule
Aritméticos & \texttt{+} \enspace \texttt{-} \enspace \texttt{*} \enspace \texttt{/} \enspace \texttt{\%} (módulo, só inteiros) e o \texttt{-} unário \\
Bit a bit & \texttt{\&} \enspace \texttt{|} \enspace \texttt{\textasciicircum} \enspace \texttt{\textasciitilde} (inversão) \\
Deslocamentos & \texttt{<<} \enspace \texttt{>>} \enspace \texttt{>>>} (à direita, preservando o sinal) \\
Relacionais & \texttt{<} \enspace \texttt{>} \enspace \texttt{<=} \enspace \texttt{>=} \enspace \texttt{==} \enspace \texttt{!=} \\
Lógicos & \texttt{\&\&} \enspace \texttt{||} \enspace \texttt{!} \\
Incremento & \texttt{++} (como comando ou em expressão, inclusive em elemento de \textit{array}) \\
\bottomrule
\end{tabular}
\end{table}

As ausências importam tanto quanto as presenças.

\begin{aviso}
Não existem em \cmm{}: o decremento \texttt{-{}-}, os operadores compostos (\texttt{+=}, \texttt{-=}, \texttt{*=} e afins), o operador ternário \texttt{?:}, o \textit{cast} explícito e os ponteiros; os símbolos \texttt{*} e \texttt{\&} são apenas multiplicação e E bit a bit. Escreva \texttt{x = x - 1;} onde escreveria \texttt{x-{}-;}.
\end{aviso}

\subsection{Operadores custam hardware}

Cada operador usado pela primeira vez no programa instancia o circuito correspondente na ULA do processador gerado, e o terminal TASM anuncia cada instância durante a compilação. Divisão e módulo são os blocos mais caros; os deslocamentos, os mais baratos. Daí o idioma do exemplo condutor, \code{soma >> 2} em vez de \code{soma / 4}. A aritmética de complexos aceita as quatro operações, cada uma expandindo para o conjunto de operações reais e imaginárias necessárias.

\section{Estruturas de controle}

Todas as estruturas clássicas do C existem, com a mesma sintaxe: o \texttt{if} com \texttt{else}; o \texttt{while} e o \texttt{do while}; o \texttt{for}, que o compilador reescreve internamente como um \texttt{while} equivalente; o \texttt{switch} com \texttt{case} e \texttt{default}, incluindo o \textit{fall-through} real do C, em que a execução continua no próximo \texttt{case} na ausência de \code{break}; e \code{break}, \code{continue} e \code{return}, os dois primeiros válidos apenas dentro de laços.

\begin{lstlisting}[style=cmm]
// exemplo: saturador
int satura(int x, int lim)
{
    if (x > lim)  return lim;
    if (x < -lim) return -lim;
    return x;
}
\end{lstlisting}

\section{Funções}

As funções seguem a sintaxe do C, \texttt{tipo nome(tipo p1, tipo p2, \ldots)}, com chamadas por argumentos posicionais e retorno por \code{return}. O número de argumentos é conferido em cada chamada, e a contagem errada é erro de compilação; conversões de tipo em parâmetros e no retorno geram aviso. A profundidade de chamadas aninhadas é limitada pela pilha de sub-rotinas, a diretiva \code{\#SDEPTH}.

Duas restrições merecem destaque. \textit{Arrays} não podem ser passados como parâmetro; trabalhe com \textit{arrays} globais ou no próprio \code{main()}. E a recursão não é suportada no fluxo \cmm{}: cada variável, incluindo parâmetros e locais, vive num endereço fixo da memória de dados, de modo que uma função que chamasse a si mesma sobrescreveria os próprios dados. É justamente esse endereçamento fixo que permite à AURORA mostrar cada variável por nome nas formas de onda. Se o algoritmo é essencialmente recursivo, o caminho C++ atende (Capítulo~\ref{cap:cmm-avancado}).

\begin{lstlisting}[style=cmm]
int add(int a, int b)
{
    return a + b;
}

int triple_sum(int x, int y, int z)
{
    return x + y + z;
}

void main()
{
    while (1)
    {
        int v = in(0);
        out(0, add(v, 1));
        out(0, triple_sum(v, 1, 2));
    }
}
\end{lstlisting}

\section{Arrays}

A linguagem aceita \textit{arrays} de uma e duas dimensões, com tamanho fixo, e uma forma de inicialização peculiar e muito útil:

\begin{lstlisting}[style=cmm]
int  x[128];               // sem inicializacao
int  h[8]    "coefs.txt";  // inicializado por arquivo, em tempo de compilacao
int  m[8][8] "matriz.txt"; // 2D, tambem por arquivo
\end{lstlisting}

A inicialização por arquivo lê os valores de um arquivo de texto ao lado do fonte e os grava diretamente na imagem da memória de dados do processador: os coeficientes do seu filtro já nascem na memória, sem custo de inicialização em tempo de execução.

Há duas formas de indexação. A normal, \code{x[k]}, e a bit-reversa, \code{x[k)}, que se distingue pelo fecha-parênteses: o índice é lido com os bits invertidos, o endereçamento clássico da FFT (\textit{Fast Fourier Transform}), com a quantidade de bits invertidos definida pela diretiva \code{\#FFTSIZ} (Capítulo~\ref{cap:cmm-avancado}). Elementos de \textit{array} aceitam o incremento, como em \code{hist[k]++;}.

\section{Variáveis e escopo}

Declarações podem aparecer no topo do programa ou dentro de funções e blocos. Na prática, porém, todo nome vive num endereço fixo e único da memória de dados, e declarar duas variáveis com o mesmo nome é erro. Pense nas variáveis \cmm{} como registradores nomeados do seu circuito: é exatamente isso que elas se tornam, e é por isso que cada uma aparece por nome na forma de onda da simulação.

% !TEX root = ../main.tex
\chapter{\cmm{}: entrada, saída e biblioteca padrão}
\label{cap:cmm-io-stdlib}

\section{Entrada e saída}
\label{sec:io}

O processador conversa com o mundo por portas numeradas, criadas pelas diretivas \code{\#NUIOIN} e \code{\#NUIOOU}, e todo o tráfego passa por quatro funções intrínsecas: \code{in(p)} lê um inteiro da porta de entrada \code{p}; \code{fin(p)} lê convertendo para \code{float}; \code{out(p, expr);} escreve o valor da expressão na porta de saída \code{p}; e \code{fout(p, expr);} escreve em formato \code{float}.

O número da porta é um literal inteiro, validado contra as diretivas: usar \code{in(2)} num processador com \code{\#NUIOIN 1} é erro de compilação. No hardware, \code{in()} é uma leitura com \textit{handshake}, na qual o processador sinaliza \code{req\_in} e aguarda o dado, e \code{out()} pulsa \code{out\_en} com o dado estável no barramento, o suficiente para acoplar FIFOs e conversores. No nosso \code{media_movel}, \code{x[0] = in(0);} lê a amostra e \code{out(0, soma >> 2);} publica a média: uma porta de cada lado basta.

\section{Funções especiais}

Cinco funções existem para economizar hardware, cada uma evitando um bloco caro da ULA ou uma sequência de instruções. A \code{norm(x)} divide pelo valor da diretiva \code{\#NUGAIN}, uma potência de dois, e entrega a normalização sem instanciar o divisor completo; funciona só com inteiros. A \code{pset(x)} devolve \code{x} quando positivo e zero quando negativo, o equivalente de \code{if (x<0) x=0;} numa única instrução. A \code{abs(x)} é o valor absoluto e, sobre um complexo, a magnitude. A \code{sign(x, y)} devolve \code{y} com o sinal de \code{x}. E a \code{copy(x, y);} copia os bits de \code{x} em \code{y} sem conversão nem checagem de tipo, uma reinterpretação crua, a ser usada com consciência.

\section{Funções não lineares e de arredondamento}

A biblioteca cobre raiz e trigonometria (\code{sqrt}, \code{sin}, \code{cos}, \code{tan}, \code{atan}), as hiperbólicas (\code{sinh}, \code{cosh}, \code{tanh}), exponencial e logaritmo (\code{exp}, \code{log}, \code{pow}) e os arredondamentos \code{floor}, \code{ceil} e \code{round}, que retornam \code{float}, com o \code{round} arredondando o meio para longe do zero.

Essas funções não são circuitos dedicados: são macros de \textit{assembly} otimizadas, como o método de Newton na \code{sqrt} e séries nas trigonométricas, injetadas no programa quando usadas. Elas custam instruções e ciclos, não blocos de ULA, exceto pelas operações de ponto flutuante que elas próprias empregam. A \code{pow(x, y)} merece nota: com expoente inteiro constante, o compilador gera a sequência mínima de multiplicações; com expoente inteiro variável, um laço em tempo de execução; nos demais casos, a identidade $x^y = e^{y \ln x}$.

\section{Funções para complexos}

Seis funções operam sobre o tipo \code{comp}: \code{real(z)} e \code{imag(z)} extraem as partes; \code{fase(z)} devolve o argumento; \code{mod2(z)} devolve a magnitude ao quadrado, $a^2+b^2$, mais barata que \code{abs} por dispensar a raiz; \code{complex(re, im)} monta um complexo a partir de dois \code{float}; e \code{conj(z)} devolve o conjugado.

\begin{lstlisting}[style=cmm]
comp x; comp y; comp r;
x = 1.0 + 2.0i;
y = 3.0 + 4.0i;
r = x * y;
fout(0, real(r));
fout(0, imag(r));
\end{lstlisting}

\begin{aviso}
Nem tudo se aplica a \code{comp}: \code{sqrt}, \code{exp}, \code{log}, \code{pow}, \code{sign} e \code{pset} sobre complexos são erro de compilação, assim como o incremento \code{z++}. O módulo \code{\%} e a \code{norm()} são exclusivos de inteiros.
\end{aviso}

\section{Álgebra linear em notação de Dirac}
\label{sec:dirac}

O recurso mais idiossincrático da \cmm{} é a escrita de operações vetoriais e matriciais em notação \textit{bra-ket}, familiar a quem vem da física. Sendo \code{a} e \code{b} vetores (\textit{arrays}) e \code{M} uma matriz, o produto interno é a expressão \texttt{<a|b>}; o comando \texttt{a \# |0>;} zera o vetor; \texttt{a \# |M|b>;} calcula o produto matriz-vetor; \texttt{a \# c|b>;} escala um vetor por uma constante; \texttt{A \# |a><b|;} forma o produto externo; e \texttt{out(p, c|a>);} emite o vetor escalado por uma porta de saída.

Essas construções expandem para os laços e instruções indexadas adequados, e são a forma idiomática de escrever filtros FIR, correlações e projeções em \cmm{}. Um FIR direto, por exemplo, é uma linha, \texttt{y = <h|x>;}, com \code{h} guardando os coeficientes, inicializados por arquivo, e \code{x} a janela de amostras.

% !TEX root = ../main.tex
\chapter{\cmm{}: recursos avançados}
\label{cap:cmm-avancado}

\section{Interrupção: \texttt{\#PRACA}}

O processador SAPHO tem um pino de interrupção (\code{itr}). Quando ele pulsa, o contador de programa desvia para o ponto do código marcado pela diretiva \code{\#PRACA} --- que, diferentemente das diretivas de configuração, aparece \textbf{no meio do programa}, marcando a linha de retorno. O uso típico é reiniciar o laço de processamento quando o hardware externo sinaliza um novo evento (um \emph{trigger} de aquisição, por exemplo), sem esperar o programa completar a volta.

\section{Sincronização com o hardware: \texttt{\#TOAQUI}}
\label{sec:toaqui}

A diretiva \code{\#TOAQUI}, colocada num ponto do programa, faz o pino \code{cheguei} do processador pulsar sempre que a execução passa por ali. É um farol para o mundo externo: \enquote{cheguei neste ponto do programa}. Serve para sincronizar blocos de hardware com fases do algoritmo --- e é também o mecanismo que a AURORA usa internamente: no \textbf{teste de hardware} (terminal THTEST) e nas simulações com Verilator, a IDE insere um \code{\#TOAQUI} ao final do \code{main()} para detectar o término do programa e encerrar a simulação no instante certo.

\section{FFT e o endereçamento bit-reverso}

A diretiva \code{\#FFTSIZ n} configura o circuito de reversão de bits para FFTs de $2^n$ pontos, usado pela indexação \code{x[k)} (Capítulo~\ref{cap:cmm-tipos-operadores}). Na FFT radix-2, os resultados saem em ordem bit-reversa; com \code{x[k)} a reordenação é gratuita --- o hardware inverte os bits do índice na leitura/escrita, sem instruções extras. O repositório do YANC traz um projeto de exemplo completo (\code{proc\_fft}) com FFT em \cmm{}.

\section{Quanto custa cada construção}
\label{sec:custo-hardware}

O lema \enquote{pague pelo que usar} tem uma consequência prática: \textbf{o custo em FPGA do seu processador é proporcional ao subconjunto da linguagem que você usa}. Durante a compilação, o terminal TASM informa cada bloco de hardware instanciado à medida que os opcodes aparecem pela primeira vez:

\begin{itemize}
  \item chamar uma função (\code{CAL}) instancia a memória de pilha de sub-rotinas;
  \item usar arrays instancia o circuito de endereçamento indexado; \code{x[k)}, o de reversão de bits;
  \item \code{/} instancia o divisor inteiro; \code{\%}, o módulo; as versões \code{float}, os blocos de ponto flutuante correspondentes;
  \item ao final, o TASM resume o percentual da ISA e da ULA efetivamente usados pelo programa.
\end{itemize}

Regras de bolso para hardware enxuto: prefira deslocamentos a divisões por potências de 2 (ou \code{norm()}); prefira \code{mod2()} a \code{abs()} quando só compara magnitudes; evite \code{float} se a dinâmica do sinal cabe em inteiros; e lembre que a primeira ocorrência paga o bloco --- a segunda ocorrência do mesmo operador é \enquote{grátis}.

\section{O caminho C++}
\label{sec:caminho-cpp}

Além da \cmm{}, o YANC compila \textbf{C++} para o mesmo processador, por uma toolchain paralela: o pré-processador \tool{cpppp} (com \code{\#include}, \code{\#define} completo e compilação condicional) seguido do compilador \tool{cppcomp}, que aceita um subconjunto de C99 com extensões de C++ e emite o mesmo assembly intermediário.

O que o caminho C++ oferece que a \cmm{} não tem:

\begin{itemize}
  \item \textbf{ponteiros} (para indexar arrays de tamanho fixo; sem alocação dinâmica);
  \item \textbf{recursão}, com quadros de pilha de verdade;
  \item \textbf{structs} simples (POD) e até recursos de C++ como \emph{lambdas} e destrutores (RAII);
  \item cabeçalhos padrão adaptados: \file{array}, \file{cmath}, \file{cstdint}, \file{vector} e outros.
\end{itemize}

O preço: variáveis locais deixam de ter endereço fixo, e com isso \textbf{deixam de aparecer por nome nas formas de onda} --- o vínculo íntimo entre código e onda, marca do fluxo \cmm{}, se perde em parte. Os parâmetros do processador no fluxo C++ têm padrões próprios (palavra de 32 bits, mantissa 23, expoente 8), ajustáveis por \code{\#pragma yanc}.

\begin{nota}
Recomendação prática: comece todo projeto em \cmm{}. Migre para C++ apenas quando precisar especificamente de recursão, ponteiros ou structs --- e mantenha os processadores de sinal críticos em \cmm{}, onde a visibilidade de simulação é total.
\end{nota}

\section{As mensagens do compilador}
\label{sec:mensagens-compilador}

Os diagnósticos do YANC são bilíngues --- acompanham o idioma da IDE (Capítulo~\ref{cap:configuracoes}) --- e têm personalidade: um erro de sintaxe rende \enquote{Erro de sintaxe na linha N} com um comentário bem-humorado, e esquecer o \code{main()} provoca a pergunta clássica sobre onde ele está. Por trás do humor, a informação é séria; as classes principais:

\begin{longtable}{@{}lp{0.60\linewidth}@{}}
\toprule
\textbf{Classe} & \textbf{Exemplos} \\
\midrule
\endhead
Estrutura & falta de \code{main()}; \code{break}/\code{continue} fora de laço; contagem errada de parâmetros. \\
Nomes & variável inexistente; variável já declarada; uso do \code{i} reservado. \\
Tipos & \code{\%}/\code{norm()} com não-inteiro; intrínsecos não implementados para \code{comp}; \code{++} em complexo. \\
Faixas & estouro de inteiro para o \code{\#NUBITS} escolhido; \code{float} fora da faixa do formato configurado; porta de E/S inexistente. \\
Avisos & conversões implícitas entre tipos (atribuição, parâmetro, retorno); \code{float}/\code{comp} em condição ou índice. \\
\bottomrule
\end{longtable}

Erros aparecem no terminal TCMM com a linha clicável --- o clique leva o editor à linha exata (Capítulo~\ref{cap:terminais}). A Aurora Intelligence também sabe explicar qualquer mensagem do compilador (ferramenta \code{lookup\_compiler\_message}, Capítulo~\ref{cap:aurora-intelligence}).


\part{O editor}
% !TEX root = ../main.tex
\chapter{O editor}
\label{cap:editor}

O editor da AURORA é construído sobre o \tool{Monaco}, o mesmo motor de edição do Visual Studio Code. Multi-cursor, minimapa, localizar e substituir, desfazer ilimitado e a ergonomia geral que você já conhece vêm de fábrica; este capítulo cobre o que a AURORA acrescenta por cima.

\section{Abas}

Cada arquivo aberto vira uma aba na barra superior do editor, com um ponto indicando alterações não salvas. As abas podem ser reordenadas por arraste. Salva-se a aba atual com \tecla{Ctrl+S} e todas com \tecla{Ctrl+Shift+S}; fecha-se a atual com \tecla{Ctrl+W}, reabre-se a última fechada com \tecla{Ctrl+Shift+T}, e \tecla{Ctrl+N} cria um arquivo novo.

\section{Editor dividido}
\label{sec:split}

O botão de divisão, que flutua no canto do painel em foco, abre o arquivo atual num novo painel lado a lado, até o limite de três painéis, cada um com a sua barra de abas. O detalhe importante é que painéis exibindo o mesmo arquivo compartilham o mesmo documento: editar de um lado atualiza o outro ao vivo, com desfazer e estado de salvamento unificados. É a forma natural de manter o \file{.cmm} à esquerda e o \file{.v} gerado à direita, ou duas regiões distantes do mesmo arquivo.

\figurapendente{Editor dividido em dois painéis}{O media\_movel.cmm à esquerda e o media\_movel.v gerado à direita.}

\section{Realce de sintaxe}

A \cmm{} tem gramática de realce própria, cobrindo palavras-chave, diretivas, funções da biblioteca, literais complexos e a notação de Dirac; os \code{\#define} do programa são recolhidos dinamicamente e coloridos como constantes. O \textit{assembly} do SAPHO também tem gramática própria. O Verilog e o SystemVerilog usam a gramática do Monaco, enriquecida com \textit{tokens} semânticos do tree-sitter, que igualmente atende C e C++ no caminho alternativo de compilação.

\section{Diagnósticos e navegação}

Dois servidores de linguagem (LSP, \textit{Language Server Protocol}) trabalham nos arquivos Verilog. O \tool{Verible} é o servidor completo: diagnósticos de sintaxe e estilo enquanto se digita, formatação, \textit{outline} de símbolos, \textit{hover}, ir à definição e referências. O \tool{slang} complementa com a análise semântica do projeto inteiro, elaborando todos os módulos em conjunto para apontar erros de elaboração e portas incompatíveis, além de oferecer autocompletar; por ser mais pesado, pode ser ligado e desligado com \tecla{Ctrl+Alt+S}. Os problemas aparecem sublinhados no código, e tudo funciona em regime de melhor esforço: sem um servidor disponível, a edição segue sem ele.

\section{Formatação automática}

O atalho \tecla{Shift+Alt+F} formata o documento atual, para C, C++ e \cmm{}, com o \tool{clang-format} empacotado. Se o projeto tiver um arquivo \file{.clang-format} próprio, ele é respeitado.

\section{Busca no projeto}

O atalho \tecla{Ctrl+Shift+F}, ou o botão de busca no cabeçalho da árvore, abre o modal Buscar nos Arquivos: termo com opções de diferenciar maiúsculas, palavra inteira e expressão regular, resultados agrupados por arquivo e clique levando à linha. A busca ignora o \file{.git} e as pastas de artefatos.

\section{Seleção e IA}

Ao selecionar um trecho de código, um pequeno botão em forma de estrela aparece junto à seleção: ele envia o trecho diretamente à Aurora Intelligence, já contextualizado. É o caminho mais curto para pedir uma explicação ou perguntar por que algo não compila (Capítulo~\ref{cap:aurora-intelligence}).

\section{Conflitos com o disco}

Se um arquivo aberto for alterado por outro programa, a AURORA detecta a divergência e pergunta o que fazer no diálogo Arquivo Modificado Externamente, oferecendo manter a versão do editor, usar a versão do disco ou salvar e recarregar. Nada é sobrescrito sem a sua decisão.

\section{Atalhos de teclado}

Os atalhos citados neste capítulo, todos reconfiguráveis nas Configurações, estão reunidos com os demais no Apêndice~\ref{ap:atalhos}. Dentro do editor valem ainda os atalhos usuais do Monaco, como \tecla{Ctrl+F} para localizar, \tecla{Alt+Click} para múltiplos cursores e \tecla{Ctrl+/} para comentar a linha. A posição do cursor é exibida na barra de status.


\part{Compilação e simulação}
% !TEX root = ../main.tex
\chapter{Compilação}
\label{cap:compilacao}

Com o \code{media_movel.cmm} escrito, é hora de transformá-lo em hardware. Este capítulo explica o que acontece quando cada botão de compilação é acionado e quais arquivos aparecem no projeto.

\section{O pipeline do YANC em três estágios}
\label{sec:pipeline-yanc}

Todo clique em \botao{Compilar \cmm{}} executa, por processador, a cadeia de três compiladores do YANC. O \tool{cmmcomp} traduz o \file{<proc>.cmm} em \textit{assembly} simbólico, o \file{<proc>.asm}, junto com os registros e as tabelas que ligam cada instrução à linha do fonte. O \tool{appcomp} faz a primeira passada sobre o \textit{assembly}, coletando os parâmetros do processador e resolvendo os endereços de variáveis e rótulos. O \tool{asmcomp}, na segunda passada, gera o processador em Verilog (\file{Hardware/<proc>.v}), as duas imagens de memória, \file{<proc>_inst.mif} com o programa e \file{<proc>_data.mif} com os dados, e o \textit{testbench} (\file{Simulation/<proc>_tb.v}).

O terminal TCMM mostra a saída do primeiro estágio; o TASM, a dos outros dois, incluindo os avisos de recurso instanciado que revelam o custo de hardware do código (Seção~\ref{sec:custo-hardware}).

\begin{info}
O processador gerado instancia os moldes de HDL do SAPHO que acompanham a instalação: o núcleo, a ULA, as memórias. O programa não é um arquivo carregado depois; as imagens \file{.mif} nascem embutidas no projeto de hardware, e reprogramar significa recompilar e sintetizar de novo.
\end{info}

\section{Os botões de compilação}

O botão \botao{Compilar \cmm{}} roda o pipeline completo da seção anterior para os processadores do projeto, do fonte ao Verilog, memórias e \textit{testbench}. O \botao{Sintetizar Verilog} verifica os arquivos Verilog com o Icarus Verilog, apontando erros de sintaxe e elaboração no terminal TVERI, com as linhas clicáveis. O \botao{Analisar Verilog} percorre o fluxo completo até as formas de onda: compila o que for preciso, roda a simulação e abre o visualizador (Capítulo~\ref{cap:simulacao}). O \botao{Fast Sim} simula sem gravar ondas, disponível com o Verilator selecionado, para quando só interessam as saídas em arquivo. O \botao{PRISM} sintetiza com o Yosys e abre o diagrama de RTL (Capítulo~\ref{cap:prism}), exigindo um \textit{top-level} definido. O \botao{Teste do processador sintetizado} constrói um \textit{harness} C++ com o Verilator em torno do \file{<proc>.v} e exercita apenas as portas de entrada e saída, o teste de fumaça do processador como caixa-preta, com progresso no THTEST. E o \botao{Cancelar} interrompe todos os processos de compilação e simulação em andamento, sem fechar a IDE.

O botão \botao{Iniciar Compilação}, o raio da barra de status, é um atalho para o fluxo principal, e todos esses comandos também existem na paleta (\tecla{Ctrl+Shift+P}), no grupo de compilação.

Antes de compilar, a AURORA salva os arquivos automaticamente. O fluxo de ondas exige um \textit{testbench}, papel que o \file{<proc>_tb.v} gerado cumpre, e o PRISM exige um \textit{top-level}; a barra de status avisa o que falta, e o processador ativo acompanha o \file{.cmm} em foco no editor.

\section{Onde cada artefato aparece}

Depois de compilar o \code{media_movel}, o projeto ganha esta forma:

\begin{lstlisting}
media_movel/
|-- Software/
|   |-- media_movel.cmm        (seu fonte)
|   |-- media_movel.asm        (assembly gerado)
|   `-- pc_media_movel_mem.txt (tabela PC -> linha, p/ ondas)
|-- Hardware/
|   |-- media_movel.v          (o processador em Verilog)
|   |-- media_movel_inst.mif   (imagem da memoria de programa)
|   `-- media_movel_data.mif   (imagem da memoria de dados)
`-- Simulation/
    `-- media_movel_tb.v       (testbench gerado)
\end{lstlisting}

O \file{.v} e os \file{.mif} da pasta \file{Hardware} são exatamente o que se leva ao Quartus ou ao Vivado na hora de sintetizar para o FPGA real. O \textit{testbench} gerado produz \textit{clock} e \textit{reset}, alimenta o processador com os estímulos de \file{input_<i>.txt}, um valor por linha e um arquivo por porta de entrada, e grava as saídas em \file{output_<i>.txt}, além do despejo de ondas.

\begin{nota}
Segurança de execução: a AURORA só executa binários da própria cadeia de ferramentas, validados contra uma lista fechada, e os botões não passam por um \textit{shell}. É por isso que a IDE não oferece a execução de comandos arbitrários fora do terminal TCMD.
\end{nota}

% !TEX root = ../main.tex
\chapter{Os terminais}
\label{cap:terminais}

A área inferior da janela reúne seis terminais em abas. Cinco são \textbf{consoles de saída} --- cada etapa da toolchain escreve no seu, então você sempre sabe onde procurar --- e um é um \textbf{shell interativo} de verdade.

\figurapendente{Área de terminais com o TCMM ativo após uma compilação}{Mostrar as seis abas, alguns cartões de mensagem (sucesso/aviso) e os filtros do topo.}

\section{Quem escreve onde}

\begin{longtable}{@{}llp{0.52\linewidth}@{}}
\toprule
\textbf{Aba} & \textbf{Nome} & \textbf{Conteúdo} \\
\midrule
\endhead
\textbf{TCMM} & Terminal \cmm{} & Saída do compilador \cmm{} (\tool{cmmcomp}): erros, avisos e progresso da tradução para assembly. \\
\textbf{TASM} & Terminal ASM & Saída do montador (\tool{appcomp}/\tool{asmcomp}): geração do Verilog, memórias, e os avisos de \enquote{recurso instanciado}. \\
\textbf{TVERI} & Terminal Verilog & Saída do Icarus Verilog e da síntese (Yosys/PRISM): diagnósticos de Verilog. \\
\textbf{TWAVE} & Terminal Wave & Execução da simulação (vvp, Verilator, cocotb) e abertura dos visualizadores de onda. \\
\textbf{THTEST} & Terminal de teste de hardware & O teste de E/S do processador sintetizado, com barra de progresso própria. \\
\textbf{TCMD} & Terminal de shell & Um \textbf{PowerShell interativo} completo (Seção~\ref{sec:tcmd}). \\
\bottomrule
\end{longtable}

Quando você dispara uma ação na toolbar, a AURORA troca automaticamente para a aba certa conforme as etapas avançam.

\section{Cartões, filtros e o modo verboso}

As mensagens dos consoles chegam como \textbf{cartões} classificados: \textcolor{statusError}{\textbf{erros}}, \textcolor{statusWarning}{\textbf{avisos}}, \textcolor{statusSuccess}{\textbf{sucessos}} e dicas/informações --- o mesmo código de cores da IDE. No topo da área ficam:

\begin{itemize}
  \item os \textbf{filtros por tipo} com contadores (clique para mostrar/ocultar cada classe); eles atuam quando o \textbf{Modo Verboso} está desligado nas Configurações --- com o modo ligado (padrão), tudo aparece;
  \item \botao{Limpar} --- esvazia o terminal atual;
  \item \botao{Exportar log} --- salva o conteúdo em arquivo;
  \item \botao{Recarregar aplicação} --- reinicia a interface da IDE.
\end{itemize}

Cada terminal retém até 5000 entradas.

\section{Números de linha clicáveis}

Erros que citam arquivo e linha --- \code{arquivo.v:42:} do Icarus/Yosys ou \enquote{linha 17} do \cmm{} --- viram \textbf{links}: o clique abre o arquivo no editor, na linha exata. É o caminho mais rápido do erro à correção.

\section{TCMD: o shell integrado}
\label{sec:tcmd}

A aba \textbf{TCMD} é um terminal PowerShell real (xterm), rodando na máquina com um \emph{prompt} próprio da AURORA:

\begin{itemize}
  \item digite qualquer comando --- \code{git}, \code{python}, \code{cd}, scripts do projeto --- e tecle \tecla{Enter};
  \item o diretório atual e as variáveis de ambiente \textbf{persistem} entre comandos, como num terminal comum;
  \item o menu de contexto da árvore de arquivos tem \menu{Abrir no Terminal Integrado}, que abre o TCMD já na pasta escolhida;
  \item é também o shell que a Aurora Intelligence usa quando você autoriza a ferramenta \code{run\_in\_terminal} (Capítulo~\ref{cap:aurora-intelligence}).
\end{itemize}

\begin{nota}
Use o TCMD para tarefas auxiliares (Git avançado, scripts Python de análise dos \file{output_*.txt}, utilitários). Para compilar e simular, prefira os botões da toolbar --- eles cuidam de argumentos, caminhos e ordem das etapas por você.
\end{nota}

% !TEX root = ../main.tex
\chapter{Simulação}
\label{cap:simulacao}

Simular é executar o processador gerado, ciclo a ciclo, num simulador de Verilog --- vendo cada variável, cada instrução e cada porta no tempo. A AURORA orquestra tudo com um clique em \botao{Analisar Verilog}; este capítulo explica as escolhas disponíveis e como controlá-las.

\section{Dois motores: Icarus Verilog e Verilator}
\label{sec:motores}

O seletor de simulador na toolbar alterna o motor global:

\begin{longtable}{@{}lp{0.36\linewidth}p{0.38\linewidth}@{}}
\toprule
 & \textbf{Icarus Verilog} (padrão) & \textbf{Verilator} \\
\midrule
\endhead
Como funciona & Interpreta o Verilog (\tool{iverilog} compila, \tool{vvp} executa). & Converte o design em C++ e compila um executável nativo. \\
Velocidade & Boa para designs pequenos/médios. & Cerca de 5--10$\times$ mais rápido em simulações longas. \\
Visibilidade & \textbf{Expõe os sinais internos} do processador (variáveis, acumulador, PC\ldots). & Apenas os sinais do topo do \emph{testbench}. \\
Uso típico & Depuração: ver o programa rodando por dentro. & Execuções longas e testes de regressão. \\
\bottomrule
\end{longtable}

Ao alternar, uma notificação resume a troca. Para o dia a dia de depuração do \code{media_movel}, fique no Icarus; quando quiser rodar milhões de ciclos, mude para o Verilator.

\section{Dois tipos de testbench: Verilog e cocotb}

O \emph{testbench} marca quem estimula o processador:

\begin{itemize}
  \item \textbf{Verilog} (\file{_tb.v}) --- o gerado pela compilação, ou um seu. O gerado cria clock e reset, alimenta cada porta de entrada com o conteúdo de \file{input_<i>.txt} (um valor por linha) e grava cada porta de saída em \file{output_<i>.txt}.
  \item \textbf{Python/cocotb} (\file{.py}) --- marque um arquivo Python como \emph{testbench} (menu de contexto da árvore, que também oferece \menu{Novo testbench cocotb (.py)}) e a AURORA roda o \emph{framework} cocotb com o Python embutido, sobre o motor selecionado. Ideal para estímulos complexos, checagens automáticas e comparação com modelos NumPy.
\end{itemize}

Os dois tipos cruzam com os dois motores: são \textbf{quatro caminhos} que convergem no mesmo arquivo de ondas (FST). A simulação roda uma única vez por clique.

\section{Parâmetros da simulação}

No popover \botao{Configurações de simulação do processador} (Seção~\ref{sec:proc-config}): frequência de clock (MHz), número de ciclos e a exibição de arrays nas ondas. O tempo simulado estimado é mostrado ali mesmo. Se a saída do seu programa não aparece, a causa mais comum é simulação curta demais --- aumente o número de ciclos.

\section{Escolhendo os sinais das ondas}
\label{sec:wave-config}

Por padrão, o \emph{dump} grava os sinais do escopo do \emph{testbench}. O botão \botao{Configuração de ondas} abre o modal \textbf{Sinais para Dump}, onde você escolhe exatamente o que gravar:

\begin{itemize}
  \item lista de todos os sinais descobertos no projeto, com busca (\tecla{Ctrl+F} dentro do modal), opções de maiúsculas/minúsculas e expressão regular;
  \item \botao{Selecionar tudo}, \botao{Limpar}, \botao{Restaurar padrão};
  \item filtro \menu{Mostrar apenas sinais do processador};
  \item a opção \menu{Surfer: manter janelas abertas} (para comparar execuções lado a lado);
  \item rodapé com o contador de selecionados e \botao{Salvar}.
\end{itemize}

\figurapendente{Modal Configuração de Ondas}{Com alguns sinais do \code{media\_movel} filtrados e selecionados.}

Menos sinais = arquivos menores e simulação mais leve --- relevante em execuções longas.

\subsection{Quem manda no dump: a ordem de precedência}

Quando a simulação vai gravar ondas, a AURORA decide a lista de sinais nesta ordem:

\begin{enumerate}
  \item um \textbf{\file{.gtkw} ativo} no seletor da toolbar --- os sinais referenciados nele mandam;
  \item a sua seleção no modal \textbf{Configuração de ondas};
  \item um \code{\$dumpvars} escrito à mão no seu \emph{testbench} --- respeitado, nada é injetado;
  \item o padrão: todo o escopo do \emph{testbench}.
\end{enumerate}

O seu \emph{testbench} nunca é modificado: a instrumentação acontece numa cópia temporária.

\section{O seletor de \texttt{.gtkw}}

Ao lado do botão de ondas, o seletor de \file{.gtkw} escolhe o \textbf{layout ativo}: \enquote{padrão} usa o layout curado gerado pela AURORA; \menu{+ Adicionar arquivo .gtkw\ldots} registra um layout seu (salvo de dentro do GTKWave). O layout ativo também dirige a precedência do \emph{dump}, acima.

\section{Rodando}

Clique em \botao{Analisar Verilog}. Acompanhe no terminal TWAVE: construção, execução (\enquote{Executando simulação VVP\ldots} ou a compilação do Verilator) e a abertura do visualizador escolhido. O Capítulo~\ref{cap:ondas} assume daqui.

Para execuções sem ondas, use \botao{Fast Sim} (Verilator): o programa roda na velocidade máxima e as saídas ficam nos \file{output_<i>.txt}.

\section{Teste de hardware (THTEST)}

O botão \botao{Teste do processador sintetizado} valida o processador como caixa-preta: um \emph{harness} C++ do Verilator alimenta as portas de entrada e confere as saídas, sem \emph{dump} de ondas, com barra de progresso no THTEST. Para detectar o fim do programa, a AURORA usa o mecanismo \code{\#TOAQUI}/\code{cheguei} (Seção~\ref{sec:toaqui}): o marcador é inserido automaticamente ao final do \code{main()} quando necessário.

\section{Simulando o \texttt{media\_movel}}

Roteiro completo do exemplo condutor:
\begin{enumerate}
  \item crie \file{Simulation/input_0.txt} com uma rampa e um degrau (um inteiro por linha);
  \item \botao{Compilar \cmm{}} --- confira no TCMM/TASM que tudo passou;
  \item confira o motor (Icarus) e o visualizador (GTKWave) nos seletores;
  \item \botao{Analisar Verilog} --- o GTKWave abre com o layout curado;
  \item observe \code{x[0]}\ldots\code{x[3]}, \code{soma} e a porta de saída suavizando o degrau; a trilha \cmm{} mostra qual linha do programa executa em cada ciclo;
  \item os valores de saída também estão em \file{output_0.txt}, prontos para conferir num script.
\end{enumerate}


\part{Visualização}
% !TEX root = ../main.tex
\chapter{Formas de onda: GTKWave e Surfer}
\label{cap:ondas}

Depois que a simulação gera o arquivo de ondas, a AURORA abre um visualizador para inspecionar cada sinal do processador ciclo a ciclo. Há dois visualizadores disponíveis. O \tool{GTKWave} é o padrão e já vem empacotado: um \textit{fork} próprio do laboratório, o \repo{gtkwave-nipscern}, com tema escuro e ajustes de usabilidade. O \tool{Surfer} é um visualizador moderno escrito em Rust, opcional, do qual a AURORA usa o \textit{fork} \texttt{surfer-aurora}, mantido pelo grupo no GitLab com decodificadores específicos do SAPHO. Ambos leem o mesmo formato de despejo, o FST, de modo que trocar de visualizador não exige simular de novo.

\figurapendente{Formas de onda de uma simulação do processador media\_movel no GTKWave}{Mostrar as faixas de clock, as trilhas de Assembly e C± em sincronia, as variáveis e as portas.}

\section{O que a AURORA prepara para você}
\label{sec:curadoria-ondas}

Abrir um despejo bruto num visualizador qualquer significa começar do zero, sem nenhum sinal selecionado e tudo em binário. A AURORA elimina esse trabalho gerando, a cada simulação, um \textit{layout} curado. Cada processador do projeto vira uma seção própria, com divisores nomeados e, no Surfer, um grupo colapsável. Os sinais de topo, \textit{clock}, \textit{reset} e interrupção, vêm primeiro; seguem as portas de entrada e saída em amarelo, as instruções em violeta, as variáveis em laranja e as \textit{flags} ao final. Variáveis do tipo \code{comp} são decodificadas para a forma legível $a+bi$ com o conversor \tool{comp2gtkw} do YANC.

O destaque são as trilhas de texto que mostram, a cada ciclo de \textit{clock}, o mnemônico \textit{assembly} executado e a linha \cmm{} correspondente, avançando em sincronia com o circuito: é o seu programa rodando dentro do hardware, visível linha a linha. Essas trilhas vêm dos arquivos de tradução gerados na compilação. No GTKWave, toda essa curadoria chega como um arquivo \file{.gtkw}; no Surfer, como um arquivo de estado \file{.surf.ron}. Nos dois casos, o visualizador já abre pronto para leitura.

\section{Escolhendo o visualizador}

O seletor de visualizador fica na \textit{toolbar}, ao lado do seletor de simulador, e vale para o aplicativo inteiro, lembrado entre sessões. Se o Surfer estiver selecionado mas o executável não estiver presente, a AURORA recua para o GTKWave sem erro.

\section{GTKWave}
\label{sec:gtkwave}

O GTKWave empacotado não é o original puro. O \textit{fork} do laboratório traz tema escuro por padrão, com o \textit{canvas} na paleta da AURORA, ajuste automático de \textit{zoom} ao abrir, de modo que a onda inteira caiba na tela, rótulos justificados à esquerda, ícones modernizados e \textit{zoom} animado.

No uso cotidiano, os sinais são adicionados a partir da árvore de escopos à esquerda; o cursor primário se posiciona com o clique e o secundário com o botão do meio, com a barra exibindo a diferença de tempo entre eles; o formato de exibição de cada sinal muda pelo menu de contexto, em \menu{Data Format}; o \textit{zoom} responde aos botões de lupa e a \tecla{Ctrl} com a roda do mouse; e a seleção atual de sinais, cores e ordem pode ser salva em um \file{.gtkw} próprio por \menu{File}, \menu{Write Save File}. Após uma nova simulação, \menu{File}, \menu{Reload Waveform} recarrega o arquivo.

\begin{nota}
Se você salvar o seu próprio \file{.gtkw} no projeto e o marcar como ativo no seletor da AURORA, ele passa a ter precedência sobre o \textit{layout} automático: a simulação seguinte grava exatamente os sinais que o seu arquivo referencia (Capítulo~\ref{cap:simulacao}).
\end{nota}

\section{Surfer}
\label{sec:surfer}

O Surfer é um visualizador com interface acelerada por GPU e carregamento preguiçoso dos arquivos, confortável em simulações grandes. A AURORA baixa o binário do \textit{fork} \texttt{surfer-aurora} automaticamente, com verificação de integridade, e o lança com o \textit{layout} curado.

Entre os recursos que valem conhecer estão os formatos de exibição ricos, de binário e hexadecimal a ponto fixo e ponto flutuante IEEE em vários tamanhos; os grupos colapsáveis, um por processador no \textit{layout} da AURORA, essenciais em projetos multiprocessados; as trilhas de \textit{assembly} e \cmm{}, instaladas automaticamente como tradutores de mapeamento; os cursores e a navegação rápidos, com uma linha de comando interna de autocompletar difuso; e o desenho analógico de sinais, útil para ver amostras de sinais processados como curvas. Os complexos do SAPHO também funcionam: a AURORA pré-decodifica os sinais \code{comp} e instala o resultado como tradutor, de modo que a faixa já abre na forma $a+bi$.

Uma limitação conhecida da versão empacotada no Windows é o recarregamento automático do arquivo de ondas, que não dispara; por isso a AURORA fecha e reabre a janela do Surfer a cada nova simulação. No modal de configuração de ondas há uma opção de manter as janelas abertas, para comparar execuções lado a lado.

% !TEX root = ../main.tex
\chapter{PRISM: o visualizador de RTL}
\label{cap:prism}

O \tool{PRISM} (\emph{Processor Rendering Interface for Schematic Models}) desenha o seu projeto como um \textbf{diagrama de circuito navegável}: módulos, portas e conexões, do top-level até dentro do processador. É a resposta visual à pergunta \enquote{que hardware o meu código virou?}.

\figurapendente{Janela do PRISM com o diagrama do projeto}{O top-level do \code{MeuFiltro} renderizado; se possível, um segundo print já dentro do módulo do processador.}

\section{Como abrir}

Botão \botao{PRISM} na toolbar (requer um \textbf{top-level} definido no projeto --- Seção~\ref{sec:arquivos-do-projeto}). A AURORA então:

\begin{enumerate}
  \item compila o que estiver pendente e verifica a sintaxe dos Verilog;
  \item reúne todos os fontes do projeto (HDL do SAPHO, sintetizáveis, top-level, \file{Hardware/} de cada processador);
  \item sintetiza com o \tool{Yosys} em modo diagrama (sem mapeamento tecnológico --- o desenho preserva a estrutura RTL do seu código);
  \item renderiza cada módulo em SVG (via \tool{netlistsvg}) e abre a janela do PRISM.
\end{enumerate}

A saída do Yosys corre no terminal TVERI. A síntese também alimenta a \textbf{visão Hierarquia} da árvore de arquivos (Seção~\ref{sec:tres-arvores}).

\section{Navegando no diagrama}

A janela do PRISM é independente da principal (1400$\times$900, maximizada ao abrir), com o mesmo tema escuro:

\begin{itemize}
  \item \textbf{pan e zoom} com o mouse;
  \item \textbf{clique} num módulo entra nele --- o diagrama do submódulo já está materializado, a navegação é instantânea;
  \item \emph{breadcrumbs} no topo mostram o caminho na hierarquia e permitem voltar;
  \item \textbf{duplo clique} numa célula abre o \textbf{código-fonte} correspondente no editor principal, na linha exata da declaração --- do desenho ao código num gesto;
  \item \botao{Recompile} refaz a síntese e o diagrama após mudanças.
\end{itemize}

Os blocos do processador SAPHO usam \emph{skins} próprias --- símbolos desenhados para a ULA, memórias, pilhas --- em vez de retângulos genéricos.

\section{Modo simulação (DigitalJS)}

Além do diagrama estático, o PRISM tem um modo \textbf{Simular}: o circuito é convertido para o \tool{DigitalJS} e vira um esquemático \emph{interativo}, onde você alterna entradas e vê os valores se propagarem pelas conexões, ao vivo. É uma lupa didática para circuitos pequenos:

\begin{itemize}
  \item o limite é de \textbf{3000 células} --- acima disso o PRISM recusa e sugere a visão esquemática comum;
  \item a síntese para este modo tem \emph{timeout} de 45 segundos.
\end{itemize}

\begin{nota}
Use o PRISM cedo e com frequência: compilar o \code{media_movel} e abrir o diagrama depois de cada mudança de código torna palpável a relação entre construções \cmm{} e custo de hardware --- adicione uma divisão ao programa e veja o divisor aparecer no desenho.
\end{nota}


\part{Aurora Intelligence}
% !TEX root = ../main.tex
\chapter{Aurora Intelligence}
\label{cap:aurora-intelligence}

A \tool{Aurora Intelligence} é a assistente de IA integrada à AURORA. Ela conversa sobre o seu projeto --- \cmm{}, Verilog, arquitetura SAPHO, mensagens do compilador --- e, com a sua permissão, \textbf{age} sobre a IDE: edita arquivos, compila, dispara simulações, seleciona sinais de onda, opera o Git. Este capítulo cobre a configuração e o uso.

\begin{info}
A Aurora Intelligence funciona no modelo \emph{traga a sua chave} (BYOK): você conecta a sua própria conta de um provedor de IA. O provedor e os modelos próprios do NIPS-CERN estão no roteiro do laboratório, mas ainda não foram treinados --- até lá, a assistente opera sobre provedores externos.
\end{info}

\section{Abrindo o painel}

Clique no botão \botao{Assistente IA} (estrela) na toolbar. O painel abre à direita, empurrando o editor; a largura é ajustável e lembrada. No cabeçalho: o ícone do provedor ativo, \botao{Chat history} (conversas anteriores), \botao{New chat} e \botao{Close}.

\figurapendente{Painel da Aurora Intelligence com uma conversa}{Uma pergunta sobre o \code{media\_movel} respondida, mostrando um cartão de ferramenta executada e o composer embaixo.}

\section{Configurando um provedor}
\label{sec:ai-config}

Em \botao{Configurações do Aurora} $\to$ aba \textbf{Assistente IA}, há um cartão por provedor:

\begin{longtable}{@{}llp{0.40\linewidth}@{}}
\toprule
\textbf{Provedor} & \textbf{Chave} & \textbf{Observações} \\
\midrule
\endhead
OpenAI & \texttt{sk-proj-\ldots} & Chave em \emph{platform.openai.com}. \\
Anthropic (Claude) & \texttt{sk-ant-api03-\ldots} & Chave em \emph{console.anthropic.com}. \\
Google (Gemini) & \texttt{AIza\ldots} & Chave em \emph{aistudio.google.com}. \\
DeepSeek & \texttt{sk-\ldots} & Chave em \emph{platform.deepseek.com}. \\
Groq & \texttt{gsk\_\ldots} & Chave em \emph{console.groq.com}. \\
Ollama (local) & --- & Sem chave: aponta para o Ollama rodando na sua máquina (\code{http://localhost:11434/v1}); o botão \botao{Detect} encontra a instância e os modelos. \\
\bottomrule
\end{longtable}

Cada cartão tem os botões \botao{Salvar}, \botao{Testar} (responde \enquote{Conectado --- N ms}), \botao{Obter uma chave} (abre o console do provedor) e \botao{Remover chave}, além do selo \textbf{Configurada}/\textbf{Sem chave} e do campo \textbf{Modelo} para escolher um modelo específico (ou manter o padrão do provedor).

\begin{nota}
\textbf{Segurança das chaves:} as chaves são cifradas pelo cofre de credenciais do Windows (DPAPI) e usadas somente pelo processo principal da AURORA --- nunca circulam pela interface nem aparecem em logs. A IDE só exibe \emph{se} há chave, nunca o valor.
\end{nota}

\subsection{CLIs por assinatura}

Alternativa às chaves de API: se você assina o \tool{Claude Code} (Anthropic) ou o \tool{Codex} (OpenAI), a Aurora Intelligence pode usar as CLIs oficiais com o login da sua assinatura. A AURORA baixa a CLI sob demanda (com verificação de integridade), a autenticação é o OAuth da própria ferramenta --- a IDE nunca vê o seu token --- e as respostas chegam no mesmo painel de chat. Nesse modo, as ferramentas da AURORA são expostas à CLI por um servidor MCP local e privado.

\section{O composer}

Na caixa de mensagem (\enquote{Pergunte à Aurora Intelligence\ldots}):

\begin{itemize}
  \item seletores de \textbf{Provedor} e \textbf{Modelo}, e \textbf{esforço de raciocínio} (para modelos que suportam);
  \item \botao{Anexar arquivos ou imagens} --- anexos multimodais;
  \item \textbf{Uso} --- consumo da conversa;
  \item \textbf{Permissões} --- o modo de autorização de ferramentas (abaixo);
  \item \tecla{Enter} envia; \botao{Parar geração} interrompe.
\end{itemize}

Selecionar código no editor e clicar na estrela do widget de seleção (Seção~\ref{cap:editor}) envia o trecho já contextualizado.

\section{Ferramentas: o que a assistente pode fazer}
\label{sec:ai-tools}

A Aurora Intelligence age por \textbf{ferramentas curadas} --- as mesmas operações dos botões da IDE, nada além. Elas se dividem em dois níveis:

\begin{itemize}
  \item \textbf{Leitura} (executam sem confirmação): ler o arquivo ativo e os abertos, a árvore do projeto, a saída dos terminais, a lista de processadores, sinais de onda, configurações; consultar a referência embutida de diretivas, opcodes e mensagens do compilador; status/log/diff do Git.
  \item \textbf{Escrita} (pedem confirmação): editar/criar/renomear/excluir arquivos, salvar, abrir arquivos e splits, compilar (\code{compile\_all}/\code{compile\_step}), cancelar compilação, criar processadores e projetos, mudar configurações, selecionar sinais de onda, abrir o Surfer, rodar comandos no TCMD, e as mutações de Git (stage, commit, push\ldots).
\end{itemize}

Cada ação de escrita aparece como um \textbf{cartão de confirmação} no chat --- \botao{Permitir} ou \botao{Negar} --- antes de executar. O seletor \textbf{Permissões} do composer ajusta o rigor: confirmar tudo, confirmar apenas escritas (padrão) ou permitir tudo.

\begin{perigo}
O modo \enquote{permitir tudo} deixa a assistente editar arquivos e rodar comandos sem perguntar. Use-o apenas em projetos versionados no Git, onde qualquer alteração é revisável e reversível (Capítulo~\ref{cap:git}).
\end{perigo}

\section{Conversas}

Cada conversa é salva localmente; o \botao{Chat history} reabre conversas antigas com contexto completo. Sem conexão, o painel indica \enquote{A Aurora Intelligence está offline}.

\section{Para que ela é boa (no fluxo SAPHO)}

\begin{itemize}
  \item \emph{\enquote{Por que este \cmm{} não compila?}} --- ela lê o terminal, conhece as mensagens do YANC e propõe a correção;
  \item \emph{\enquote{Quanto hardware custa esta função?}} --- ela consulta a referência de opcodes e o \file{.asm} gerado;
  \item \emph{\enquote{Escreva um testbench cocotb que compare com um modelo NumPy}};
  \item \emph{\enquote{Selecione os sinais das portas e abra as ondas}} --- ferramentas de onda por comando de voz\ldots{} escrito.
\end{itemize}


\part{Ferramentas de apoio}
% !TEX root = ../main.tex
\chapter{Dagr: controle de versão}
\label{cap:git}

\tool{Dagr} (do nórdico antigo \enquote{dia}; o painel usa a runa \emph{dagaz} como marca) é o painel de controle de versão da AURORA: uma interface completa de \textbf{Git} sobre o repositório do projeto aberto, com integração ao GitHub. Você não precisa instalar nem conhecer a linha de comando do Git para versionar seus projetos SAPHO --- embora o TCMD esteja ali para quem preferir.

\section{Abrindo}

Botão \botao{Controle de Versão} na toolbar (ícone de ramificação, com um contador de alterações pendentes) ou o indicador do GitHub na barra de status. Abre o modal \textbf{Dagr --- Controle de Versão}.

\figurapendente{Modal do Dagr com alterações pendentes}{A aba Changes com arquivos preparados/não preparados, o diff de um deles e a caixa de commit.}

Se o projeto ainda não é um repositório, o Dagr oferece inicializá-lo (\code{git init}) --- e criar um \file{.gitignore} adequado está a um clique no menu de contexto da árvore.

\section{A aba Changes}

O fluxo cotidiano:

\begin{enumerate}
  \item a lista de \textbf{alterações} mostra cada arquivo modificado/criado/excluído; o clique abre o \emph{diff} colorido (adições/remoções) embutido;
  \item \textbf{prepare} o que vai entrar no commit --- por arquivo ou \botao{Preparar tudo} (o inverso, \botao{Despreparar}, também existe);
  \item escreva o \textbf{Resumo do commit} (e a descrição opcional) e clique em \botao{Commit}.
\end{enumerate}

Apoios: \botao{Amend} (emenda o último commit), \botao{Desfazer último} (desfaz o commit mantendo as alterações), descartar alterações de um arquivo, e avisos para arquivos binários ou \emph{diffs} muito grandes.

A árvore de arquivos também participa: cada arquivo ganha a \textbf{decoração de status Git} (modificado, novo, ignorado) na cor correspondente.

\section{A aba History}

O histórico de commits do ramo atual: autor, data, mensagem; o clique mostra os arquivos e o \emph{diff} de cada commit.

\section{Ramos e stashes}

\begin{itemize}
  \item \textbf{Branches}: criar, trocar (com a opção \botao{Guardar e trocar}, que faz \emph{stash} automático das alterações pendentes) e mesclar.
  \item \textbf{Stashes}: a prateleira de alterações guardadas, com \botao{Restaurar} e \botao{Descartar}.
\end{itemize}

\section{GitHub: publicar, clonar, sincronizar}

\begin{itemize}
  \item \botao{Entrar com GitHub} --- dois métodos: o \emph{device flow} (\enquote{Abrimos o GitHub no seu navegador. Digite este código:}) ou um token clássico com escopo \code{repo} (o Dagr traz um guia embutido de como criar o token). O token é cifrado no cofre do sistema; \botao{Desconectar} o remove.
  \item \botao{Publicar} --- cria o repositório no GitHub a partir do projeto local e envia tudo.
  \item \botao{Push} / \botao{Pull} / \botao{Fetch} --- sincronização com o remoto, com progresso ao vivo (o \emph{pull} usa \emph{autostash}: suas alterações locais são guardadas e reaplicadas).
  \item \botao{Clonar repositórios} --- traz um repositório remoto para o disco; se contiver um \file{.spf}, o Dagr oferece \botao{Abrir projeto no SAPHO}.
\end{itemize}

\begin{nota}
Versione desde o primeiro dia. O \file{.spf} é JSON amigável a \emph{diffs}, os fontes \cmm{}/Verilog são texto --- e um repositório é o pré-requisito para usar a Aurora Intelligence com tranquilidade (todo passo dela vira um \emph{diff} revisável no Dagr).
\end{nota}

A Aurora Intelligence expõe o mesmo backend como ferramentas (\code{git\_status}, \code{git\_commit}, \code{git\_push}\ldots), com as mutações sujeitas a confirmação --- Capítulo~\ref{cap:aurora-intelligence}.

% !TEX root = ../main.tex
\chapter{Configurações}
\label{cap:configuracoes}

O botão \botao{Configurações do Aurora} na toolbar abre o modal de configurações, organizado em sete seções na navegação lateral. O rodapé traz \botao{Restaurar Padrões} e \botao{Salvar Alterações}.

\figurapendente{Modal de Configurações}{A seção Geral aberta, mostrando a navegação lateral com as sete seções.}

\section{Geral}

\begin{itemize}
  \item \textbf{Habilitar Tooltips} (ligado por padrão) --- as dicas ricas ao passar o mouse pelos botões da IDE.
  \item \textbf{Confiar em links externos} (desligado por padrão) --- quando ligado, links do chat da Aurora Intelligence abrem no navegador sem confirmação.
\end{itemize}

\section{Aparência}

Troca do \textbf{ícone da aplicação}: envie uma imagem sua ou \botao{Restaurar ícone padrão}.

\begin{info}
A AURORA tem um único tema visual --- o tema escuro aurora, cuidadosamente construído --- e ele não é configurável. É uma decisão de design da plataforma.
\end{info}

\section{Idioma}

\textbf{Português} (padrão) ou \textbf{English}. A escolha muda a interface \emph{e} as mensagens do compilador YANC (que recebe a flag \code{-pt}/\code{-en} correspondente) --- erros de compilação chegam no idioma da IDE.

\section{Terminal}

\textbf{Modo Verboso} (ligado por padrão): mostra todas as mensagens de compilação. Desligue para ativar os filtros por classe (erros, avisos, sucessos, dicas) nos terminais --- útil quando só interessam os problemas.

\section{Atalhos de Teclado}

Todos os atalhos da IDE, editáveis: clique num atalho, tecle a nova combinação (esc cancela). Combinações em conflito são recusadas com o aviso \enquote{Este atalho já está em uso}. A lista padrão está no Apêndice~\ref{ap:atalhos}.

\section{Assistente IA}

Os cartões de provedores da Aurora Intelligence --- chaves, modelos, testes de conexão (Seção~\ref{sec:ai-config}).

\section{Sobre}

A ficha da plataforma: versão instalada, a equipe (desenvolvimento e orientação), os vínculos institucionais (UFJF, NIPS-CERN) e agradecimentos (CAPES, CNPq, FAPEMIG), os links oficiais e o catálogo de \textbf{software de terceiros} com as respectivas licenças --- do Icarus Verilog ao Monaco, transparência completa sobre o que compõe a AURORA.


\part{Manutenção da plataforma}
% !TEX root = ../main.tex
\chapter{Atualizações automáticas}
\label{cap:atualizacoes}

O SAPHO evolui rápido, e a AURORA cuida de se manter atual: um atualizador automático embutido verifica, baixa e instala novas versões --- sempre com o seu aval.

\section{Como funciona, do seu ponto de vista}

\begin{enumerate}
  \item \textbf{Verificação silenciosa}: cerca de seis segundos após abrir a IDE, a AURORA consulta o canal oficial de \emph{releases} (o repositório \repo{sapho} no GitHub). Se você já está na última versão, nada acontece --- nenhum diálogo.
  \item \textbf{Oferta}: havendo versão nova, abre a \textbf{janela de atualização}: versão nova e atual, o tamanho do download e o \emph{changelog} bilíngue com o que mudou. Nada é baixado sem você pedir.
  \item \textbf{Download}: ao aceitar, a barra de progresso mostra percentual, MB, velocidade e tempo restante. Você pode continuar trabalhando.
  \item \textbf{Instalação}: com o download pronto, o botão vira \botao{Reiniciar e instalar}. A AURORA fecha, o instalador roda e a IDE reabre na versão nova.
  \item \textbf{Confirmação}: na primeira abertura pós-atualização, um \emph{toast} discreto confirma: \enquote{O SAPHO foi atualizado para a versão X.Y.Z}.
\end{enumerate}

\figurapendente{Janela de atualização}{Com uma versão nova disponível: changelog, tamanho do download e o botão de baixar. Se possível, um segundo print com o download em progresso.}

A integridade do pacote é verificada (\emph{hash} SHA-512 publicado junto com a \emph{release}) antes de instalar; downloads são incrementais quando possível. Em caso de problema --- sem internet, sem espaço em disco, sem permissão --- a janela explica em linguagem clara e a IDE continua funcionando na versão atual.

\begin{info}
\textbf{Dois repositórios, por desenho:} o desenvolvimento da AURORA acontece no repositório \repo{aurora}; as versões estáveis são publicadas no repositório \repo{sapho}, que é o canal que o seu instalador e o atualizador consomem. Você sempre recebe releases estáveis, nunca o estado intermediário do desenvolvimento.
\end{info}

\section{A toolchain acompanha}

Cada instalador carrega, embutida, a versão da toolchain testada com aquela versão da IDE --- os compiladores do YANC, os simuladores, o GTKWave, os moldes de HDL. Atualizar a AURORA atualiza o conjunto todo, em versões que foram validadas juntas. Não há nada para você gerenciar manualmente.

\section{Este manual também se atualiza}

O manual que você lê é parte do ecossistema versionado: cada mudança relevante na AURORA, no YANC ou nos componentes reflete numa edição nova. A versão do manual acompanha a da plataforma.

% !TEX root = ../main.tex
\chapter{Solução de problemas}
\label{cap:solucao-problemas}

Um capítulo de pronto-socorro: os tropeços mais comuns e o caminho de volta.

\section{Compilação e simulação}

\begin{longtable}{@{}p{0.40\linewidth}p{0.54\linewidth}@{}}
\toprule
\textbf{Sintoma} & \textbf{Causa e solução} \\
\midrule
\endhead
\enquote{Erro de sintaxe na linha N} no TCMM & Erro de \cmm{}. Clique na linha para ir ao ponto; confira as ausências da linguagem (sem \code{--}, sem \code{+=}, sem ponteiros --- Capítulo~\ref{cap:cmm-tipos-operadores}). \\
\enquote{yo, where's the main()?} & Todo programa \cmm{} precisa de \code{void main()}. \\
Erro citando porta de E/S inexistente & O índice em \code{in()}/\code{out()} excede \code{\#NUIOIN}/\code{\#NUIOOU}. Aumente a diretiva (ou o campo no Hub) e recompile. \\
Total de bits recusado no Hub & A regra é \texttt{NUBITS} = \texttt{NBMANT} + \texttt{NBEXPO} + 1, e o Ganho deve ser potência de 2. \\
Botão de ondas não faz nada / aviso na barra de status & Falta \emph{testbench} (ou top-level, para o PRISM). Compile o \cmm{} para gerar o \file{_tb.v}, ou marque o seu pelo menu de contexto. \\
A simulação termina antes dos resultados & Poucos ciclos. Aumente \textbf{Number of clocks} no popover de configurações do processador. \\
Simulação lenta demais & Troque o motor para Verilator (5--10$\times$ mais rápido; só sinais de topo) ou reduza os sinais no \botao{Configuração de ondas}. \\
Sinais internos sumiram das ondas & Você está no Verilator, que só expõe o topo do \emph{testbench}. Volte ao Icarus para depurar por dentro. \\
Processo travado & \botao{Cancelar} na toolbar interrompe toda a toolchain sem fechar a IDE. \\
\enquote{Binário Icarus Verilog não encontrado} & A pasta \file{components/} da instalação está incompleta ou foi movida. Reinstale o SAPHO pelo instalador oficial. \\
\bottomrule
\end{longtable}

\section{Projeto e arquivos}

\begin{longtable}{@{}p{0.40\linewidth}p{0.54\linewidth}@{}}
\toprule
\textbf{Sintoma} & \textbf{Causa e solução} \\
\midrule
\endhead
\enquote{Projeto Não Encontrado} & O \file{.spf} foi movido ou apagado. Abra pelo novo caminho; o item sai dos recentes sozinho. \\
Aviso de arquivos faltantes & Arquivos registrados no \file{.spf} sumiram do disco. Restaure-os ou dispense o aviso para limpar as referências. \\
\enquote{Arquivo Modificado Externamente} & Outro programa alterou um arquivo aberto. Escolha entre manter a versão do editor, a do disco, ou salvar e recarregar. \\
Arrastar arquivo para a árvore não funciona & A importação por arraste aceita \file{.v}, \file{.sv}, \file{.vh} e \file{.py} (cocotb). Arquivos \file{.gtkw} entram pelo seletor da toolbar. \\
\bottomrule
\end{longtable}

\section{Visualizadores}

\begin{itemize}
  \item \textbf{Selecionei o Surfer e abriu o GTKWave}: o executável do Surfer não está presente --- a AURORA degrada para o GTKWave de propósito. Reinstale/atualize para obter o componente.
  \item \textbf{O Surfer não recarregou a onda nova}: comportamento esperado no Windows --- a AURORA fecha e reabre a janela a cada simulação.
  \item \textbf{O PRISM recusou o modo Simular}: o circuito passa de 3000 células; use o diagrama esquemático comum.
\end{itemize}

\section{Aurora Intelligence}

\begin{itemize}
  \item \textbf{\enquote{A Aurora Intelligence está offline}}: sem conexão com o provedor. Verifique a rede e teste a chave na aba Assistente IA (\botao{Testar}).
  \item \textbf{Modelo indisponível}: a AURORA recua automaticamente para o modelo padrão do provedor; escolha outro no seletor do composer.
  \item \textbf{Ollama não detectado}: confirme que o serviço local está rodando e use \botao{Detect} no cartão do Ollama.
\end{itemize}

\section{Onde estão os logs}

Para relatar um problema, anexe o log principal: \file{\%APPDATA\%\textbackslash Aurora-IDE\textbackslash logs\textbackslash main.log} (o anterior fica em \file{main.old.log}). O conteúdo dos terminais pode ser salvo com \botao{Exportar log}. Issues são bem-vindas nos repositórios do \texttt{nipscernlab} no GitHub.


\appendix
\part{Apêndices}
% !TEX root = ../main.tex
\chapter{Referência de diretivas}
\label{ap:diretivas}

\section{Diretivas de configuração}

Uma por linha, no topo do \file{.cmm}, sem ponto e vírgula. A coluna do padrão indica o valor assumido pela cadeia de ferramentas quando a diretiva é omitida; o Hub de Processadores da AURORA escreve todas explicitamente.

\begin{longtable}{lllp{0.42\linewidth}}
\toprule
Diretiva & Argumento & Padrão & Efeito \\
\midrule
\endhead
\code{\#PRNAME} & nome & & Nome do processador; deve casar com o nome no projeto e nos artefatos. \\
\code{\#NUBITS} & inteiro & 23 & Largura da palavra da ULA e dos inteiros, em complemento de dois; deve valer a soma de mantissa, expoente e um. \\
\code{\#NBMANT} & inteiro & 16 & Bits de mantissa do ponto flutuante próprio. \\
\code{\#NBEXPO} & inteiro & 6 & Bits de expoente, em complemento de dois. \\
\code{\#NDSTAC} & inteiro & 10 & Profundidade da pilha de dados. \\
\code{\#SDEPTH} & inteiro & 10 & Profundidade da pilha de sub-rotinas. \\
\code{\#NUIOIN} & inteiro & 1 & Número de portas de entrada; \code{in(p)} exige \texttt{p} menor que esse valor. \\
\code{\#NUIOOU} & inteiro & 1 & Número de portas de saída; idem para \texttt{out}. \\
\code{\#NUGAIN} & potência de 2 & 64 & Divisor fixo da função \code{norm()}. \\
\code{\#FFTSIZ} & inteiro & 8 & Bits invertidos no endereçamento bit-reverso \code{x[k)}, para FFT de $2^n$ pontos. \\
\bottomrule
\end{longtable}

Com uma única porta a ligação é direta; com mais de uma, o hardware instancia automaticamente um decodificador de endereços.

\section{Diretivas comportamentais}

Aparecem no corpo do programa, marcando pontos do código. A \code{\#PRACA} marca o ponto para onde o processador desvia quando o pino de interrupção pulsa. A \code{\#TOAQUI} marca um endereço do programa cujo alcance faz pulsar o pino \code{cheguei}; é também o mecanismo que a AURORA usa para detectar o fim do programa em testes e simulações.

\section{Pré-processador}

A construção \code{\#define NOME corpo} define uma constante simbólica, apenas na forma de objeto, sem argumentos, com o limite de 256 por programa. Não há \code{\#include} nem \code{\#ifdef} no fluxo \cmm{}.

% !TEX root = ../main.tex
\chapter{Referência da biblioteca padrão}
\label{ap:stdlib}

Todas as funções intrínsecas da \cmm{}, agrupadas. \enquote{Tipos} indica os tipos aceitos; erros de tipo são apontados na compilação.

\section{Entrada e saída}

\begin{longtable}{@{}llp{0.52\linewidth}@{}}
\toprule
\textbf{Função} & \textbf{Tipos} & \textbf{Descrição} \\
\midrule
\endhead
\code{in(p)} & \code{int} & Lê a porta de entrada \code{p}. \\
\code{fin(p)} & \code{float} & Lê a porta \code{p} convertendo para \emph{float}. \\
\code{out(p, x);} & qualquer & Escreve \code{x} na porta de saída \code{p}. \\
\code{fout(p, x);} & \code{float} & Escreve em formato \emph{float}. \\
\bottomrule
\end{longtable}

\section{Funções especiais (economia de hardware)}

\begin{longtable}{@{}llp{0.52\linewidth}@{}}
\toprule
\textbf{Função} & \textbf{Tipos} & \textbf{Descrição} \\
\midrule
\endhead
\code{norm(x)} & \code{int} & $x\,/\,$\texttt{\#NUGAIN} sem instanciar divisor. \\
\code{pset(x)} & \code{int}, \code{float} & $x$ se $x \geq 0$; senão $0$. \\
\code{abs(x)} & todos & Valor absoluto; magnitude para \code{comp}. \\
\code{sign(x, y)} & \code{int}, \code{float} & $y$ com o sinal de $x$. \\
\code{copy(x, y);} & todos & Cópia bit a bit, sem checagem de tipo. \\
\bottomrule
\end{longtable}

\section{Não-lineares (macros de assembly)}

\begin{longtable}{@{}lp{0.64\linewidth}@{}}
\toprule
\textbf{Função} & \textbf{Descrição} \\
\midrule
\endhead
\code{sqrt(x)} & Raiz quadrada (Newton--Raphson). Não aceita \code{comp}. \\
\code{sin(x)}, \code{cos(x)}, \code{tan(x)} & Trigonométricas. \\
\code{atan(x)} & Arco-tangente. \\
\code{sinh(x)}, \code{cosh(x)}, \code{tanh(x)} & Hiperbólicas. \\
\code{exp(x)}, \code{log(x)} & Exponencial e logaritmo natural. Não aceitam \code{comp}. \\
\code{pow(x, y)} & Potência; estratégia depende do expoente (constante inteiro $\to$ multiplicações; variável inteiro $\to$ laço; senão $e^{y\ln x}$). \\
\code{floor(x)}, \code{ceil(x)}, \code{round(x)} & Arredondamentos (retornam \code{float}; \code{round} afasta do zero no meio). \\
\bottomrule
\end{longtable}

\section{Complexos}

\begin{longtable}{@{}lp{0.64\linewidth}@{}}
\toprule
\textbf{Função} & \textbf{Descrição} \\
\midrule
\endhead
\code{real(z)}, \code{imag(z)} & Partes real e imaginária. \\
\code{fase(z)} & Argumento (ângulo). \\
\code{mod2(z)} & Magnitude ao quadrado. \\
\code{complex(re, im)} & Constrói um \code{comp}. \\
\code{conj(z)} & Conjugado. \\
\bottomrule
\end{longtable}

\section{Notação de Dirac}

\begin{longtable}{@{}lp{0.60\linewidth}@{}}
\toprule
\textbf{Sintaxe} & \textbf{Operação} \\
\midrule
\endhead
\texttt{<a|b>} & Produto interno (expressão). \\
\texttt{a \# |0>;} & Zera o vetor. \\
\texttt{a \# |M|b>;} & Matriz vezes vetor. \\
\texttt{a \# c|b>;} & Vetor escalado. \\
\texttt{A \# |a><b|;} & Produto externo. \\
\texttt{out(p, c|a>);} & Emite o vetor escalado pela porta \code{p}. \\
\bottomrule
\end{longtable}

\section{Famílias de instruções do processador}

Para leitura das trilhas de assembly nas ondas: os mnemônicos seguem convenções de prefixo/sufixo. \code{F\_} indica a versão em ponto flutuante; \code{S\_}/\code{SF\_}, operando vindo da pilha; \code{P\_}, empilha antes; \code{\_M}, opera sobre memória; \code{\_V}, variante vetorial/indexada.

\begin{longtable}{@{}lp{0.62\linewidth}@{}}
\toprule
\textbf{Família} & \textbf{Mnemônicos principais} \\
\midrule
\endhead
Carga/armazenamento & \code{LOD}, \code{SET}, \code{LDI}/\code{STI} (indexados), \code{ILI}/\code{ISI} (bit-reverso) \\
Pilha & \code{PSH}, \code{POP} \\
E/S & \code{INN}, \code{F\_INN}, \code{OUT} \\
Fluxo & \code{JMP}, \code{JIZ}, \code{CAL}, \code{RET}, \code{NOP} \\
Aritmética inteira & \code{ADD}, \code{MLT}, \code{DIV}, \code{MOD}, \code{NEG}, \code{ABS}, \code{SGN}, \code{PST}, \code{NRM} \\
Aritmética float & \code{F\_ADD}, \code{F\_SU1}/\code{F\_SU2}, \code{F\_MLT}, \code{F\_DIV}, \code{F\_NEG}, \code{F\_ABS}, \code{F\_ROT}, \code{F\_SCL}, \code{XPO} \\
Conversões & \code{I2F}, \code{F2I} \\
Lógica e bits & \code{AND}, \code{ORR}, \code{XOR}, \code{INV}, \code{LAN}, \code{LOR}, \code{LIN} \\
Comparações & \code{LES}, \code{GRE}, \code{EQU} (e as versões \code{F\_}) \\
Deslocamentos & \code{SHL}, \code{SHR}, \code{SRS} (equivale ao operador \code{>>>}) \\
\bottomrule
\end{longtable}

Cada mnemônico usado pela primeira vez no programa liga o bloco de hardware correspondente no processador gerado --- a lista completa e o percentual de uso da ISA aparecem no terminal TASM a cada compilação.

% !TEX root = ../main.tex
\chapter{Atalhos de teclado}
\label{ap:atalhos}

Padrões de fábrica, todos reconfiguráveis nas Configurações, na seção Atalhos de Teclado.

\begin{longtable}{lp{0.52\linewidth}}
\toprule
Atalho & Ação \\
\midrule
\endhead
\tecla{Ctrl+N} & Novo arquivo \\
\tecla{Ctrl+S} & Salvar \\
\tecla{Ctrl+Shift+S} & Salvar todos \\
\tecla{Ctrl+W} & Fechar aba \\
\tecla{Ctrl+Shift+T} & Reabrir última aba fechada \\
\tecla{Ctrl+Shift+F} & Buscar nos arquivos do projeto \\
\tecla{Ctrl+Shift+P} & Paleta de comandos \\
\tecla{Ctrl+Alt+S} & Ligar ou desligar a análise semântica (slang) \\
\tecla{Shift+Alt+F} & Formatar documento (clang-format) \\
\bottomrule
\end{longtable}

Dentro do editor valem ainda os atalhos usuais do Monaco, herdados do VS Code: \tecla{Ctrl+F} para localizar, \tecla{Ctrl+H} para substituir, \tecla{Ctrl+Z} e \tecla{Ctrl+Y} para desfazer e refazer, \tecla{Alt+Click} para múltiplos cursores e \tecla{Ctrl+/} para comentar a linha, entre outros.

% !TEX root = ../main.tex
\chapter{Glossário}
\label{ap:glossario}

\begin{description}[style=nextline,font=\bfseries]
  \item[SAPHO] \emph{Scalable-Architecture Processor for Hardware Optimization}: o processador \emph{soft-core} e, por extensão, a plataforma completa (AURORA + YANC + toolchain) e o canal de distribuição.
  \item[AURORA] \emph{Advanced Utility Running Optimized Resource Architectures}: a IDE de desktop da plataforma.
  \item[YANC] \emph{Yet Another Compiler}: a suíte de compiladores que traduz \cmm{}/C++ em processador Verilog.
  \item[\cmm{} (CMM)] A linguagem de programação da plataforma; dialeto de C com complexos nativos e notação de Dirac. Arquivos \file{.cmm}.
  \item[PRISM] \emph{Processor Rendering Interface for Schematic Models}: o visualizador de RTL da AURORA.
  \item[Aurora Intelligence] A assistente de IA integrada à AURORA.
  \item[Dagr] O painel de controle de versão (Git) da AURORA.
  \item[NIPS-CERN] O Núcleo de Instrumentação e Processamento de Sinais da UFJF, em colaboração com o CERN --- o laboratório que desenvolve a plataforma.
  \item[soft-processor / soft-core] Processador descrito em HDL e implementado na malha reconfigurável de um FPGA.
  \item[FPGA] \emph{Field-Programmable Gate Array}: circuito integrado reconfigurável.
  \item[Verilog / HDL] Linguagem de descrição de hardware usada pela plataforma.
  \item[\file{.spf}] \emph{SAPHO Project File}: o arquivo JSON que define um projeto.
  \item[\file{.asm}] O assembly intermediário do SAPHO, entre o compilador e o montador.
  \item[\file{.mif}] \emph{Memory Initialization File}: imagem de memória (programa ou dados) embutida no design.
  \item[testbench] O banco de testes da simulação: um módulo Verilog (\file{_tb.v}) ou um script Python (cocotb) que estimula o processador.
  \item[top-level] O módulo Verilog que integra os processadores e blocos do projeto; topo da síntese.
  \item[ULA] Unidade Lógico-Aritmética: o conjunto de operadores do processador; no SAPHO, só os usados são instanciados.
  \item[acumulador (ACC)] O registrador central do processador SAPHO, destino de toda operação da ULA.
  \item[Harvard] Arquitetura com memórias de programa e dados separadas.
  \item[opcode / ISA] Código de operação de uma instrução / o conjunto de instruções do processador.
  \item[Icarus Verilog / vvp] O simulador de Verilog padrão (interpreta o design).
  \item[Verilator] Simulador de alto desempenho que converte o design em executável C++.
  \item[cocotb] \emph{Framework} Python para escrever \emph{testbenches}.
  \item[Yosys] Ferramenta de síntese usada pelo PRISM e pela visão Hierarquia.
  \item[GTKWave / Surfer] Os dois visualizadores de formas de onda da plataforma.
  \item[VCD / FST] Formatos de arquivo de ondas gerados pela simulação (textual / binário compacto).
  \item[\file{.gtkw} / \file{.surf.ron}] Arquivos de layout de ondas (GTKWave / Surfer): quais sinais, cores e grupos mostrar.
  \item[Monaco] O motor de edição de código (o mesmo do VS Code).
  \item[LSP] \emph{Language Server Protocol}: o mecanismo dos diagnósticos e navegação de código (Verible, slang).
  \item[BYOK] \emph{Bring Your Own Key}: o modelo de chaves da Aurora Intelligence.
  \item[MCP] \emph{Model Context Protocol}: como as CLIs de IA acessam as ferramentas da AURORA.
  \item[bit-reverso] Endereçamento com bits do índice invertidos (\code{x[k)}), usado em FFT.
  \item[DLP] Dispositivos Lógicos Programáveis: a disciplina da UFJF que usa a plataforma no ensino.
\end{description}


\end{document}
